% Incognita paper P007: the setting paper (kind `definitions`) of the roadmap StochasticBandits.
% Written from templates/paper/paper.tex (AGENTS.md, section 5; docs/writing.md, sections 2-2a).
% Two builds of one source, after `python3 scripts/certify.py bib` (papers/library.bib):
%   reader:    latexmk -pdf -interaction=nonstopmode -halt-on-error -cd papers/P007-stochastic-setting/paper.tex
%   blueprint: latexmk -pdf -interaction=nonstopmode -halt-on-error -jobname=paper-blueprint
%                -usepretex='\def\IncognitaBlueprint{}' -cd papers/P007-stochastic-setting/paper.tex
% Blueprint conventions: every node has a \label (def:, lem:, prop:, thm:, cor:, cite:); statements
% start with \uses when they use definitions; every proof starts with an exhaustive \uses; a proof
% away from its statement has \proves. The planned Lean names are in DOSSIER.md (never \lean here).
\documentclass[11pt]{article}
\usepackage[T1]{fontenc}
\usepackage[utf8]{inputenc}
\usepackage{lmodern}
\usepackage[a4paper,margin=2.5cm]{geometry}
\usepackage{amsmath,amssymb,mathtools}
\usepackage{booktabs}
\usepackage{algorithm}
\usepackage{algpseudocode}
\usepackage{enumitem}
\usepackage{xcolor}
\usepackage[colorlinks=true,linkcolor=blue!50!black,citecolor=blue!50!black,urlcolor=blue!50!black]{hyperref}
\usepackage[round]{natbib}
\usepackage{../../templates/paper/blueprint-macros}

% ---------------------------------------------------------------- macros (\newcommand only)
% Each symbol has one meaning in the whole paper.
\newcommand{\E}{\mathbb{E}}
\newcommand{\Prob}{\mathbb{P}}
\newcommand{\R}{\mathbb{R}}
\newcommand{\ind}{\mathbf{1}}
\newcommand{\argmax}{\operatorname*{arg\,max}}
% Navigation (house style): an anchor \dt{s:...} at each definition, one macro per symbol linking to it.
\newcommand{\dt}[1]{\hypertarget{#1}{}}
\newcommand{\sy}[2]{{\hyperlink{#1}{#2}}}
\newcommand{\term}[2]{\hyperlink{#1}{#2}}
% laws and divergences
\newcommand{\Pall}{\sy{s:Pall}{\mathcal{P}(\R)}}
\newcommand{\Pone}{\sy{s:P1}{\mathcal{P}_{1}}}
\newcommand{\mean}{\sy{s:mean}{m}}
\newcommand{\KL}{\sy{s:KL}{\mathrm{KL}}}
\newcommand{\kl}{\sy{s:kl}{\mathrm{kl}}}
\newcommand{\Ber}{\sy{s:Ber}{\mathrm{Ber}}}
\newcommand{\Kinf}{\sy{s:kinf}{\mathcal{K}_{\inf}}}
% models
\newcommand{\Punit}{\sy{s:Punit}{\mathcal{P}[0,1]}}
\newcommand{\Bclass}{\sy{s:Bclass}{\mathcal{B}}}
\newcommand{\Semi}{\sy{s:S}{\mathcal{S}}}
\newcommand{\Lmom}{\sy{s:L}{\mathcal{L}}}
\newcommand{\var}{\sy{s:var}{\sigma^{2}}}
\newcommand{\kurt}{\sy{s:kurt}{\mathrm{kurt}}}
\newcommand{\Hk}{\sy{s:H}{\mathcal{H}}}
\newcommand{\Nor}{\sy{s:Nor}{\mathcal{N}}}
\newcommand{\Gau}{\sy{s:G}{\mathcal{G}}}
% instances and classes
\newcommand{\arms}{\sy{s:arms}{\alpha}}
\newcommand{\Karms}{\sy{s:K}{K}}
\newcommand{\marm}{\sy{s:ma}{m}}
\newcommand{\ms}{\sy{s:mstar}{m^{\star}}}
\newcommand{\gap}{\sy{s:gap}{\Delta}}
\newcommand{\Opt}{\sy{s:Opt}{\mathcal{A}^{\star}}}
\newcommand{\prodc}[1]{\sy{s:prod}{\textstyle\prod_{a\in\alpha}}#1_{a}}
\newcommand{\powc}[1]{\sy{s:prod}{#1^{\alpha}}}
% learners and runs
\newcommand{\Hs}{\sy{s:Hs}{\mathbb{H}}}
\newcommand{\pol}{\sy{s:pol}{p}}
\newcommand{\Hist}{\sy{s:Hist}{H}}
\newcommand{\Ap}{\sy{s:A}{A}}
\newcommand{\Xr}{\sy{s:X}{X}}
\newcommand{\Pru}{\sy{s:Prun}{\Prob}}
\newcommand{\Er}{\sy{s:Prun}{\E}}
\newcommand{\Np}{\sy{s:N}{N}}
\newcommand{\hR}{\sy{s:hR}{\hat{R}}}
\newcommand{\Reg}{\sy{s:R}{R}}
\newcommand{\hm}{\sy{s:hm}{\hat{m}}}
\newcommand{\taup}{\sy{s:tau}{\tau}}
\newcommand{\Yc}{\sy{s:Y}{Y}}
\newcommand{\bY}{\sy{s:bY}{\bar{Y}}}
\newcommand{\Zc}{\sy{s:Z}{Z}}
\newcommand{\Cons}{\sy{s:pib}{\pi}}
% Graves--Lai
\newcommand{\as}{\sy{s:astar}{a^{\star}}}
\newcommand{\Alt}{\sy{s:Alt}{\mathrm{Alt}}}
\newcommand{\Alto}{\sy{s:Alto}{\mathrm{Alt}_{0}}}
\newcommand{\GL}{\sy{s:GL}{C}}
\newcommand{\Kd}{\sy{s:Kd}{k}}
% UCB
\newcommand{\Uix}{\sy{s:U}{U}}
\newcommand{\ucb}{\hyperref[alg:ucb]{\texttt{UCB}}}

\title{The setting of stochastic multi-armed bandits with finitely many arms: models, learners,
consistency classes, and the Graves--Lai constant of unstructured classes}
\author{Claude Opus 5.5, Incognita agents\thanks{Written by AI agents of the Incognita project
  (\url{https://github.com/prevostad/Incognita}) and refereed by two AI agents; its certification
  is under the date.}}
\date{2026-11-16}

\begin{document}
\maketitle

\begin{abstract}
We fix the objects of problem-dependent analyses of stochastic bandits with finitely many arms:
models of reward laws (on $[0,1]$, Bernoulli, semi-bounded, with a bounded moment of order
$1+\varepsilon$, with a bounded kurtosis, Gaussian with unknown variance), the divergence
$\mathcal{K}_{\inf}$ of a model with a strict constraint on the mean, learners and their runs,
regret, the classes of learners of the lower bounds (consistent, $\gamma$-consistent, uniformly
fast and super-fast convergent, smarter than uniform, pairwise symmetric, monotonic, with
upper-tailed counts or regret), and the Graves--Lai constant $C(\nu;\mathcal{E})$ of a class
$\mathcal{E}$ of instances. The main result is the identity
$C(\nu;\prod_a\mathcal{D}_a)=\sum_{a:\Delta_a>0}\Delta_a/\mathcal{K}_{\inf}(\nu_a,m^\star;\mathcal{D}_a)$
in $[0,\infty]$, for every unstructured class and every instance of it with a unique optimal arm,
with no assumption on the models. Each definition comes with an object that satisfies it and one
that does not: for example, the upper confidence bound algorithm with a Hoeffding bonus is
uniformly super-fast convergent on laws on $[0,1]$ with the constant $13$, while a learner that
always pulls the same arm is not consistent. We also compute $\mathcal{K}_{\inf}$ on Bernoulli and
Gaussian laws, bound it on laws on $[0,1]$, show that it vanishes on all laws with a finite mean,
and prove that the rewards collected by count form an independent array in every run.
\end{abstract}

% ================================================================ main text
\section{Introduction}\label{sec:intro}

In the stochastic multi-armed bandit problem, a learner pulls one of finitely many arms at each
round and observes a reward drawn from the law of that arm. Since \citet{LR85} and \citet{BK96},
its problem-dependent performance is measured against a lower bound: when the learner knows a set
$\mathcal D$ of laws that contains the law of every arm, every consistent learner pulls each
suboptimal arm $a$ at least $(1-o(1))\log T/\mathcal K_{\inf}(\nu_a,m^\star;\mathcal D)$ times in
$T$ rounds, where
\begin{equation*}
  \mathcal K_{\inf}(P,\mu;\mathcal D)=\inf\{\mathrm{KL}(P,Q):Q\in\mathcal D,\ \text{the mean of }Q\text{ is larger than }\mu\}.
\end{equation*}
When the instance is only known to lie in a set $\mathcal E$ of instances with a structure, the
constant of the lower bound becomes the value of an optimisation problem over exploration rates,
the Graves--Lai constant $C(\nu;\mathcal E)$ \citep[Theorem~1]{CMP17}; see also \citet{GL97}.

Papers that combine results on these problems need these objects with one convention each: which
inequality on the mean defines $\mathcal K_{\inf}$ (strict or not), which values it may take, what
a learner and its run are, how the rounds are counted, which classes of learners the lower bounds
quantify over, and how the Graves--Lai problem treats infinite divergences and infinite weights.
Published papers differ on each point, and a paper that combines two results must check that
they agree. This paper fixes the conventions, proves the basic properties of the objects, and
checks each definition against an object that satisfies it and an object that does not.

The main result relates the two constants. For every unstructured class $\prod_a\mathcal D_a$
(each arm has its own model $\mathcal D_a$, and the learner knows nothing else) and every instance
$\nu$ of it with a unique optimal arm,
\begin{equation*}
  C\Big(\nu;\prod_{a}\mathcal D_a\Big)=\sum_{a:\Delta_a>0}\frac{\Delta_a}{\mathcal K_{\inf}(\nu_a,m^\star;\mathcal D_a)}\in[0,\infty]
\end{equation*}
(Theorem~\ref{thm:main}), with no assumption on the models beyond finite means: the divergences
may vanish or be infinite, and both sides are then computed with the conventions of
$[0,\infty]$. The identity shows that $C(\nu;\mathcal E)$ extends the Lai--Robbins constant, and a
result stated for every class $\mathcal E$ can be checked on its unstructured case.

\paragraph{Our contributions} We summarise our contributions below.
\begin{enumerate}[label=(\roman*),leftmargin=2.2em]
  \item \emph{Definitions.} The models of reward laws (Definitions~\ref{def:laws}
    and~\ref{def:models}), $\mathcal K_{\inf}$ (Definition~\ref{def:kinf}), instances and classes
    (Definitions~\ref{def:instance} and~\ref{def:product}), learners, runs, counts and regret
    (Definitions~\ref{def:learner}--\ref{def:regret}), the classes of learners of the lower bounds
    (Definitions~\ref{def:consistent}--\ref{def:smarter} and~\ref{def:pairwise}--\ref{def:family}),
    the alternatives and the Graves--Lai constant (Definition~\ref{def:gl}).
  \item \emph{The Graves--Lai constant of an unstructured class} (Theorem~\ref{thm:main}), from a
    reduction that removes the optimal arm from the problem (Lemma~\ref{lem:gl-reduction}); the
    constant vanishes when the law of the optimal arm identifies the instance
    (Proposition~\ref{prop:gl-zero}); with the lower bound of \citet[Theorem~1]{GMS19}, every
    consistent learner has $\liminf_T R_T/\log T\ge C(\nu;\mathcal D^\alpha)$ (Corollary~\ref{cor:lower}).
  \item \emph{Properties of $\mathcal K_{\inf}$}: monotonicity, its values on laws on $[0,1]$
    (bounds by $\mathrm{kl}$ and by a logarithm), closed forms on Bernoulli and Gaussian laws, its
    vanishing on all laws with a finite mean (Lemmas~\ref{lem:kinf-basic}--\ref{lem:kinf-gauss});
    properties of the semi-bounded, moment and kurtosis models (Lemmas~\ref{lem:semi}
    and~\ref{lem:kurt}).
  \item \emph{Witnesses and exclusions for the classes of learners.} The rewards collected by
    count are an independent array in every run (Lemma~\ref{lem:stack}); \ucb{} with a Hoeffding
    bonus is uniformly super-fast convergent on laws on $[0,1]$ with the constant $13$
    (Proposition~\ref{prop:ucb}), hence consistent, and its counts have tails in $1/(\log T)^2$
    (Proposition~\ref{prop:ucb-tails}); constant, uniform and switching learners delimit the
    classes (Lemmas~\ref{lem:constant}, \ref{lem:gms-classes} and~\ref{lem:tails}); the relations
    between the classes (Lemmas~\ref{lem:classes} and~\ref{lem:family}).
  \item \emph{Examples} with explicit constants (Section~\ref{sec:examples}).
\end{enumerate}

\paragraph{Outline} Section~\ref{sec:related} says where these models and classes come from.
Section~\ref{sec:setting} states the definitions, and Section~\ref{sec:main} the results on the
Graves--Lai constant, the main result included. Section~\ref{sec:properties} states the
properties of the models and of $\mathcal K_{\inf}$, and Section~\ref{sec:learners} those of the
classes of learners with their witnesses and exclusions. Section~\ref{sec:overview} explains how
the proofs work, Section~\ref{sec:examples} gives examples, and Section~\ref{sec:open} states the
limitations. The proofs are in the appendices.

% ----------------------------------------------------------------
\section{Related work and context}\label{sec:related}

This section says where each object comes from and which convention we keep when the sources
differ. Table~\ref{tab:sources} lists the sources definition by definition.

\paragraph{The lower bound and $\mathcal K_{\inf}$} \citet{LR85} proved the logarithmic lower bound
for one-parameter families and \citet{BK96} for general models, with the constant
$\mathcal K_{\inf}$. We use the form of \citet[Section~1.2]{GMS19}, with a strict inequality on
the mean and $\inf\emptyset:=+\infty$, which is also that of the textbook
\citet[Theorem~16.2]{LS}; their consistency is \citet[Definition~16.1]{LS}, and the lower bound is
proved for the larger class of uniformly fast convergent strategies in
\citet[Definition~1 and Theorem~1]{GMS19}. A source that uses the non-strict inequality $\mean(Q)\ge\mu$, such as
\citet[Section~2]{AJK21}, defines a function $\mathcal K^{\ge}_{\inf}$ with
$\lim_{\mu'\uparrow\mu}\mathcal K_{\inf}(P,\mu';\mathcal D)\le\mathcal K^{\ge}_{\inf}(P,\mu;\mathcal D)\le\mathcal K_{\inf}(P,\mu;\mathcal D)$,
since the constraint sets are nested; the two agree wherever ours is left-continuous in $\mu$, and a
result stated with the non-strict inequality must be translated before it is combined with ours.

\paragraph{Classes of learners} Uniform super-fast convergence and learners smarter than the
uniform one are from \citet[Definitions~6 and~2]{GMS19}, where they are the hypotheses of a
second-order lower bound and of lower bounds in the initial regime. The policies that we call
$\gamma$-consistent are the $\alpha$-consistent policies of \citet[Definition~5]{SAE13}, where the
lower bound degrades by the factor $1-\alpha$; we write $\gamma$ for the exponent because $\alpha$
is the set of arms. Pairwise symmetric and monotonic strategies are the hypotheses of the lower
bounds in the initial regime of \citet[Definitions~3 and~4, Theorems~3 and~4]{GMS19}, and learners with
upper-tailed counts or regret are those of \citet[Definition~1]{SA11}. These sources define a class on
a model $\mathcal D$, that is on the instances of
$\mathcal D^\alpha$; we define each class on an arbitrary set of instances, which contains their
definitions as the case $\mathcal E=\mathcal D^\alpha$.

\paragraph{Models of reward laws} Laws on $[0,1]$ are the setting of KL-UCB for general bounded
rewards \citep[Section~6]{CGMMS13} (their empirical KL-UCB, Section~5, is analysed for finitely
supported laws), of KL-UCB-switch \citep[Section~2]{GHMS22} and of non-parametric
Thompson sampling \citep[Theorem~5]{RH20}. Semi-bounded laws, on $(-\infty,1]$ with a finite mean,
are the model of IMED \citep[Section~1]{HT15}. Laws with $\E|X|^{1+\varepsilon}\le B$ are the model
of \citet[Section~2]{AJK21}; \citet[Theorem~1]{BCL13} assume the same bound on the raw moment of
order $1+\varepsilon$.
Laws with a bounded kurtosis are the scale-free model of \citet[Assumption~1]{L17}. Gaussian laws
with unknown mean and variance, for which $\mathcal K_{\inf}$ has the closed form that we recover in
Lemma~\ref{lem:kinf-gauss}, are studied by \citet[Section~2]{HT14} and \citet[Section~1]{CHK18}.

\paragraph{Structured classes and the Graves--Lai constant} The lower bound over a structured class,
with constraints indexed by the alternatives that leave the law of the optimal arm unchanged, goes
back to finite parameter spaces and to controlled Markov chains \citep[see also][]{ATA89,GL97}, as
recounted by \citet[Section~4]{VPG24}. \citet[Theorem~1]{CMP17} state it for a structured class
$\Theta$ with a unique optimal arm, as a corollary of the result on controlled Markov chains, with a
constant defined by a linear program over these alternatives. Their Section~4 gives the
unstructured case for parametric models, with the rate $1/d(\theta(x),\theta(x^\star))$ for each
suboptimal arm. \citet[Proposition~1]{VPG24} state the lower bound for closed convex classes of laws
on a common finite set of rewards.

\paragraph{Closest related work} Two published statements are special cases of
Theorem~\ref{thm:main}. The first is the unstructured case of \citet[Section~4]{CMP17}, for
parametric families in which the divergence is finite and continuous in the parameter. The second
is the case of separable bandits of \citet[Example~3, with Propositions~1 and~2]{VPG24}: a product
of closed convex sets of laws on a common finite set of rewards. There the optimal exploration rate
is stated to be $1/\mathrm{KL}(\nu_a,Q_a)$, with $Q_a$ a minimiser over the laws of the model of $a$
with a mean at least $m^\star$, on an arm $a$ whose model contains a law with a mean larger than
$m^\star$ (a ``deceitful'' arm), and $0$ on the other arms; the attainment of the infimum is their
Proposition~2, and their example~(9) has the constant $0$, our case $\mathcal K_{\inf}=+\infty$.
Neither statement gives Theorem~\ref{thm:main} for general models: in both settings
$\mathcal K_{\inf}$ of a suboptimal arm is positive and attained, whereas for general models $\mathcal K_{\inf}$ may be $0$, the infimum need
not be attained, and the laws may have unbounded supports (Gaussian or semi-bounded laws). Our
definition of $C(\nu;\mathcal E)$ follows the convention of the problem that we are given: all the
alternatives, and weights in $[0,\infty]$ with $0\cdot\infty:=0$. Lemma~\ref{lem:gl-reduction}
shows that it agrees with the formulation of \citet{CMP17} that keeps the law of the optimal arm,
and Theorem~\ref{thm:main} then follows by an argument arm by arm that handles every degenerate
value.

\paragraph{Gap with the literature} The objects themselves are not new, and the identity of
Theorem~\ref{thm:main} is known in the two cases above. What this paper adds is a single set of
conventions under which the objects are all defined, with complete proofs of their basic
properties, and Theorem~\ref{thm:main} for arbitrary models, including the degenerate values
($\mathcal K_{\inf}\in\{0,\infty\}$, an infimum that is not attained, best mean on the boundary of
the model, infinite weights). Lemma~\ref{lem:stack}, that the rewards collected by count form an
independent array, restates a known fact: it is a form of the optional skipping theorem
\citep[see also][]{D36}, and the equivalence of the run with the ``stack of rewards'' model asserted by
\citet[Section~4.6]{LS}; our proof is the one of its formalisation in Lean in the library
\citet{LML} (file \texttt{RewardByCountMeasure.lean}). The analysis of \ucb{} in
Proposition~\ref{prop:ucb} is that of UCB1 by \citet[Theorem~1]{ACF02}, written in our conventions
with explicit constants.

\begin{table}[t]
\centering
\caption{Where the definitions come from, and the convention we keep.}\label{tab:sources}
\small
\begin{tabular}{@{}p{0.30\textwidth}p{0.30\textwidth}p{0.32\textwidth}@{}}
\toprule
Object (node) & Source & Convention kept \\
\midrule
$\mathcal K_{\inf}$ (Def.~\ref{def:kinf}) & \citet[Section~1.2]{GMS19}; \citet[Theorem~16.2]{LS} & strict inequality on the mean, values in $[0,\infty]$ \\
consistency (Def.~\ref{def:consistent}) & \citet[Definition~16.1]{LS} & on any set of instances \\
$\gamma$-consistency (Def.~\ref{def:consistent}) & \citet[Definition~5]{SAE13} ($\alpha$-consistency) & $\gamma\in[0,1)$ \\
uniform fast, super-fast convergence (Def.~\ref{def:ufc}) & \citet[Definitions~1 and~6]{GMS19} & on any set of instances, $T\ge2$ \\
smarter than uniform (Def.~\ref{def:smarter}) & \citet[Definition~2]{GMS19} & on any set of instances \\
pairwise symmetric, monotonic (Defs.~\ref{def:pairwise}, \ref{def:monotonic}) & \citet[Definitions~3 and~4]{GMS19} & on any set of instances, every $T\ge1$ \\
upper-tailed counts, regret (Def.~\ref{def:tails}) & \citet[Definition~1]{SA11} & on any set of instances, pseudo-regret \\
Graves--Lai constant (Def.~\ref{def:gl}) & \citet[Theorem~1]{CMP17} & all alternatives, weights in $[0,\infty]$, $0\cdot\infty:=0$ \\
models (Def.~\ref{def:models}) & \citet{CGMMS13,HT15,AJK21,L17,HT14} & one set of laws per model \\
\bottomrule
\end{tabular}
\end{table}

% ----------------------------------------------------------------
\section{Problem formulation}\label{sec:setting}

In this section we state the definitions: laws and divergences, models, instances and classes,
learners and their runs, the classes of learners, and the Graves--Lai constant. The classical
objects come first; where we change a convention of the sources, we say so where it appears.

\subsection{Laws, divergences and models}

\begin{definition}[Laws, means and divergences]\label{def:laws}
\dt{s:Pall}A \emph{law} is a Borel probability measure on $\R$, and $\Pall$ is the set of laws.
\dt{s:P1}The set of laws with a finite mean is
\begin{equation*}
  \Pone:=\Big\{P\in\Pall:\int|x|\,dP(x)<\infty\Big\}.
\end{equation*}
\dt{s:mean}The mean of $P\in\Pone$ is $\mean(P):=\int x\,dP(x)$. \dt{s:KL}For two laws $P,Q$, the
Kullback--Leibler divergence is
\begin{equation*}
  \KL(P,Q):=\int\log\frac{dP}{dQ}\,dP\ \text{ if }P\ll Q,\qquad\KL(P,Q):=+\infty\ \text{ otherwise};
\end{equation*}
it lies in $[0,\infty]$ and vanishes if and only if $P=Q$. \dt{s:Ber}For $p\in[0,1]$, $\Ber(p)$ is
the Bernoulli law $(1-p)\delta_0+p\delta_1$, where $\delta_x$ is the Dirac mass at $x$.
\dt{s:kl}The Bernoulli divergence is, for $p,q\in[0,1]$,
\begin{equation*}
  \kl(p,q):=p\log\frac pq+(1-p)\log\frac{1-p}{1-q},
\end{equation*}
with $0\log(0/x):=0$ and $y\log(y/0):=+\infty$ for $y>0$; then $\kl(p,q)=\KL(\Ber(p),\Ber(q))$.
\end{definition}

\begin{definition}[$\Kinf$ of a model]\label{def:kinf}
\uses{def:laws}
A \emph{model} is a set $\mathcal D\subseteq\Pone$. For a model $\mathcal D$, a law $P$ and
$\mu\in\R$,
\begin{equation*}
  \dt{s:kinf}\Kinf(P,\mu;\mathcal D):=\inf\big\{\KL(P,Q):\ Q\in\mathcal D,\ \mean(Q)>\mu\big\}\in[0,\infty],
  \qquad\inf\emptyset:=+\infty.
\end{equation*}
\end{definition}

The inequality $\mean(Q)>\mu$ is strict. With it, $\Kinf(P,\mu;\mathcal D)=0$ whenever
$P\in\mathcal D$ and $\mean(P)>\mu$, and the lower bound constant of an arm whose mean equals the
best mean is never used. A model is often defined as any set of laws on $\R$; we restrict it to
laws with a finite mean, so that the constraint $\mean(Q)>\mu$ is defined, as in
\citet[Section~1.2]{GMS19}.

\begin{definition}[Models]\label{def:models}
\uses{def:laws}
\dt{s:Punit}The laws on $[0,1]$ are $\Punit:=\{P\in\Pall:P([0,1])=1\}$. \dt{s:Bclass}The Bernoulli
laws are $\Bclass:=\{\Ber(p):p\in[0,1]\}$. \dt{s:S}The semi-bounded laws are
\begin{equation*}
  \Semi:=\{P\in\Pone:P((-\infty,1])=1\}.
\end{equation*}
\dt{s:L}For $\varepsilon>0$ and $B>0$, the laws with a bounded moment of order $1+\varepsilon$ are
\begin{equation*}
  \Lmom_{\varepsilon,B}:=\Big\{P\in\Pall:\int|x|^{1+\varepsilon}\,dP(x)\le B\Big\}.
\end{equation*}
\dt{s:var}For $P\in\Pall$ with $\int x^2\,dP<\infty$, the variance is
$\var(P):=\int(x-\mean(P))^2\,dP(x)$. \dt{s:kurt}If moreover $\int x^4\,dP<\infty$ and $\var(P)>0$,
the kurtosis is
\begin{equation*}
  \kurt(P):=\frac{\int(x-\mean(P))^4\,dP(x)}{\var(P)^2}.
\end{equation*}
\dt{s:H}For $\kappa>1$, the laws with kurtosis at most $\kappa$ are
\begin{equation*}
  \Hk_\kappa:=\Big\{P\in\Pall:\int x^4\,dP<\infty,\ \var(P)>0,\ \kurt(P)\le\kappa\Big\}.
\end{equation*}
\dt{s:Nor}For $\mu\in\R$ and $v>0$, $\Nor(\mu,v)$ is the Gaussian law with mean $\mu$ and variance
$v$. \dt{s:G}The Gaussian laws with unknown mean and variance are
$\Gau:=\{\Nor(\mu,v):\mu\in\R,\ v>0\}$.
\end{definition}

All these sets are models, that is subsets of $\Pone$ (Lemma~\ref{lem:semi} for $\Lmom_{\varepsilon,B}$),
and
\begin{equation*}
  \Bclass\subseteq\Punit\subseteq\Semi\subseteq\Pone,\qquad
  \Punit\subseteq\Lmom_{\varepsilon,B}\ \text{ when }B\ge1.
\end{equation*}
The Gaussian model excludes the degenerate variance $v=0$, and the kurtosis model the laws with
$\var(P)=0$, on which the kurtosis is not defined.

\subsection{Instances and classes}

\begin{definition}[Arms and instances]\label{def:instance}
\uses{def:laws}
\dt{s:arms}The arms form a finite set $\arms$. \dt{s:K}The number of arms is $\Karms:=|\arms|\ge2$.
An \emph{instance} is a family $\nu=(\nu_a)_{a\in\arms}$ of laws $\nu_a\in\Pone$. \dt{s:ma}The mean
of arm $a$ is $\marm_a:=\mean(\nu_a)$. \dt{s:mstar}The best mean is $\ms:=\max_{a\in\arms}\marm_a$.
\dt{s:gap}The gap of arm $a$ is $\gap_a:=\ms-\marm_a\ge0$. \dt{s:Opt}The set of optimal arms is
$\Opt(\nu):=\{a\in\arms:\gap_a=0\}$; an arm $a$ is \emph{suboptimal} if $\gap_a>0$.
\end{definition}

When several instances appear, we write $\marm_a(\nu')$, $\ms(\nu')$, $\gap_a(\nu')$ and
$\Opt(\nu')$ for the quantities of $\nu'$; without argument they refer to $\nu$.

\begin{definition}[Classes of instances]\label{def:product}
\uses{def:instance, def:kinf}
A \emph{class} is a set $\mathcal E$ of instances, known to the learner. \dt{s:prod}For a family
$(\mathcal D_a)_{a\in\arms}$ of models, the \emph{unstructured class} is
\begin{equation*}
  \prodc{\mathcal D}:=\{\nu:\ \nu_a\in\mathcal D_a\ \text{for every }a\in\arms\},
\end{equation*}
and $\powc{\mathcal D}:=\prodc{\mathcal D}$ when $\mathcal D_a=\mathcal D$ for every $a$. A class is
also called a \emph{structured class}; the unstructured classes are the special case
$\prodc{\mathcal D}$.
\end{definition}

\subsection{Learners, runs, counts and regret}

\begin{definition}[Learners]\label{def:learner}
\uses{def:instance}
\dt{s:Hs}For $t\ge0$, the set of histories of length $t$ is $\Hs_t:=(\arms\times\R)^t$, with the
product $\sigma$-algebra ($\Hs_0$ is a point). \dt{s:pol}A \emph{learner} is a sequence
$\pi=(\pol_t)_{t\ge0}$ in which $\pol_t$ is a Markov kernel from $\Hs_t$ to $\arms$: after a history
$h\in\Hs_t$, the arm of round $t$ is drawn from the law $\pol_t(h,\cdot)$ on $\arms$. A learner is
also called \emph{anytime}: it is the same at every horizon. A \emph{horizon-dependent learner} is
a family $(\pi^T)_{T\ge1}$ of learners, $\pi^T$ being used when the horizon is $T$.
\end{definition}

Rounds are numbered $0,1,2,\dots$, and the horizon $T$ means the rounds $0,\dots,T-1$. Randomised
learners are included: the kernels carry the randomisation.

\begin{definition}[Runs]\label{def:run}
\uses{def:learner}
Let $\pi=(\pol_t)_{t\ge0}$ be a learner and $\nu$ an instance. \dt{s:A}A \emph{run} of $\pi$ on
$\nu$ is a probability space $(\Omega,\mathcal F,\Prob)$ with random variables $\Ap_t\in\arms$ and
\dt{s:X}$\Xr_t\in\R$, $t\ge0$, with the following property. Let
\begin{equation*}
  \dt{s:Hist}\Hist_t:=(\Ap_0,\Xr_0,\dots,\Ap_{t-1},\Xr_{t-1})\in\Hs_t
\end{equation*}
be the history before round $t$. For every $t\ge0$,
\begin{align*}
  &\text{the conditional law of }\Ap_t\text{ given }\Hist_t\text{ is }\pol_t(\Hist_t,\cdot),\\
  &\text{the conditional law of }\Xr_t\text{ given }(\Hist_t,\Ap_t)\text{ is }\nu_{\Ap_t}.
\end{align*}
\end{definition}

A run exists for every learner and every instance, and the law of the sequence
$(\Ap_t,\Xr_t)_{t\ge0}$ is the same in every run: it is the measure given by the Ionescu-Tulcea
theorem from $\pol_0,\nu,\pol_1,\nu,\dots$. \dt{s:Prun}We write $\Pru_{\pi,\nu}$ and $\Er_{\pi,\nu}$ for
probabilities and expectations of functions of $(\Ap_t,\Xr_t)_{t\ge0}$, which depend only on $\pi$
and $\nu$; the indices are dropped when $\pi$ and $\nu$ are clear.

\begin{definition}[Counts, empirical means and regret]\label{def:regret}
\uses{def:run, def:instance}
In a run of a learner $\pi$ on an instance $\nu$, \dt{s:N}the number of pulls of arm $a$ before
round $t$ is
\begin{equation*}
  \Np_a(t):=\sum_{u=0}^{t-1}\ind(\Ap_u=a).
\end{equation*}
\dt{s:hm}If $\Np_a(t)\ge1$, the empirical mean of arm $a$ before round $t$ is
\begin{equation*}
  \hm_a(t):=\frac{1}{\Np_a(t)}\sum_{u=0}^{t-1}\Xr_u\ind(\Ap_u=a).
\end{equation*}
\dt{s:hR}The pseudo-regret at horizon $T\ge1$ is $\hR_T:=\sum_{t=0}^{T-1}\gap_{\Ap_t}$, and
\dt{s:R}the regret is
\begin{equation*}
  \Reg_T(\pi,\nu):=\Er_{\pi,\nu}\big[\hR_T\big]=\sum_{a\in\arms}\gap_a\,\Er_{\pi,\nu}\big[\Np_a(T)\big].
\end{equation*}
The regret of a horizon-dependent learner $(\pi^T)_{T\ge1}$ at horizon $T$ is $\Reg_T(\pi^T,\nu)$.
\end{definition}

The second expression of the regret is the pathwise identity $\hR_T=\sum_a\gap_a\Np_a(T)$ in
expectation. A paper that numbers rounds from $1$ uses the same quantities: its $N_a(T)$, the
number of pulls among the first $T$ rounds, is our $\Np_a(T)$.

\subsection{Classes of learners}

\begin{definition}[Consistency and $\gamma$-consistency]\label{def:consistent}
\uses{def:regret, def:product}
Let $\mathcal E$ be a class and $\gamma\in[0,1)$. \dt{w:acons}A learner $\pi$ is
\emph{$\gamma$-consistent on $\mathcal E$} if, for every $\nu\in\mathcal E$ and every $p>\gamma$,
\begin{equation*}
  \lim_{T\to\infty}\frac{\Reg_T(\pi,\nu)}{T^p}=0.
\end{equation*}
\dt{w:cons}It is \emph{consistent on $\mathcal E$} if it is $0$-consistent. A horizon-dependent
learner $(\pi^T)_{T\ge1}$ is consistent on $\mathcal E$ if $\Reg_T(\pi^T,\nu)/T^p\to0$ for every
$\nu\in\mathcal E$ and every $p>0$.
\end{definition}

\begin{definition}[Uniform fast and super-fast convergence]\label{def:ufc}
\uses{def:regret, def:product}
Let $\mathcal E$ be a class. \dt{w:ufc}A learner $\pi$ is \emph{uniformly fast convergent on
$\mathcal E$} if, for every $\nu\in\mathcal E$, every suboptimal arm $a$ of $\nu$ and every
$p\in(0,1]$,
\begin{equation*}
  \lim_{T\to\infty}\frac{\Er_{\pi,\nu}[\Np_a(T)]}{T^p}=0.
\end{equation*}
\dt{w:usfc}For $C>0$, it is \emph{uniformly super-fast convergent on $\mathcal E$ with the constant
$C$} if, for every $\nu\in\mathcal E$, every suboptimal arm $a$ of $\nu$ and every $T\ge2$,
\begin{equation*}
  \Er_{\pi,\nu}[\Np_a(T)]\le\frac{C\log T}{\gap_a^2}.
\end{equation*}
\end{definition}

\begin{definition}[Smarter than the uniform learner]\label{def:smarter}
\uses{def:regret, def:product}
Let $\mathcal E$ be a class. \dt{w:smart}A learner $\pi$ is \emph{smarter than the uniform learner
on $\mathcal E$} if, for every $\nu\in\mathcal E$, every optimal arm $a\in\Opt(\nu)$ and every
$T\ge1$,
\begin{equation*}
  \Er_{\pi,\nu}[\Np_a(T)]\ge\frac{T}{\Karms}.
\end{equation*}
\end{definition}

\paragraph{Changes to the sources} The sources define these classes on the instances of a model
$\mathcal D^\alpha$; Definitions~\ref{def:consistent}--\ref{def:smarter} take an arbitrary class,
which is needed for classes that are not unstructured and contains the sources as the case
$\mathcal E=\mathcal D^\alpha$. \citet[Definition~5]{SAE13} call the property $\alpha$-consistency and take $\alpha\ge0$; we write
$\gamma$ for the exponent, since $\alpha$ is the set of arms, and take $\gamma<1$, since every
learner is $\gamma$-consistent for $\gamma\ge1$ on a class of instances with bounded gaps. Models are
subsets of $\Pone$ (Definition~\ref{def:kinf}), not arbitrary sets of laws.

\subsection{Alternatives and the Graves--Lai constant}

\begin{definition}[Alternatives and Graves--Lai constant]\label{def:gl}
\uses{def:product, def:laws}
Let $\mathcal E$ be a class and $\nu$ an instance with a unique optimal arm, \dt{s:astar}written
$\as$, so that $\Opt(\nu)=\{\as\}$. \dt{s:Alt}The alternatives to $\nu$ in $\mathcal E$ are the
instances of $\mathcal E$ in which $\as$ is not optimal:
\begin{equation*}
  \Alt(\nu;\mathcal E):=\big\{\nu'\in\mathcal E:\ \as\notin\Opt(\nu')\big\}.
\end{equation*}
\dt{s:GL}The Graves--Lai constant is
\begin{equation*}
  \GL(\nu;\mathcal E):=\inf\Big\{\sum_{a\in\arms}w_a\gap_a:\ w\in[0,\infty]^{\arms},\
  \sum_{a\in\arms}w_a\KL(\nu_a,\nu'_a)\ge1\ \text{for every }\nu'\in\Alt(\nu;\mathcal E)\Big\}\in[0,\infty],
\end{equation*}
with the conventions $0\cdot\infty:=0$ and $x+\infty:=\infty$ of $[0,\infty]$. A weight vector $w$
that satisfies every constraint is called \emph{feasible}.
\end{definition}

\paragraph{The conventions of $[0,\infty]$} A weight $w_a=\infty$ is allowed. It costs nothing on
the optimal arm, since $\gap_{\as}=0$ and $\infty\cdot0=0$, and it satisfies every constraint of an
alternative that changes the law of the optimal arm, since then $\KL(\nu_{\as},\nu'_{\as})>0$. In
quotients, $x/0:=\infty$ for $x>0$ and $x/\infty:=0$ for $x<\infty$. Several optimal arms are not
covered by Definition~\ref{def:gl}; they need a definition of their own.

% ----------------------------------------------------------------
\section{The Graves--Lai constant of an unstructured class}\label{sec:main}

In this section we first show that the optimal arm can be removed from the Graves--Lai problem,
then prove the main result, the identity on unstructured classes, and the two consequences that
check the definition: the constant vanishes when the optimal arm identifies the instance, and it
is the constant of the lower bound of \citet{GMS19} on $\powc{\mathcal D}$.

\begin{lemma}[Removing the optimal arm]\label{lem:gl-reduction}
\uses{def:gl}
Let $\mathcal E$ be a class and $\nu$ an instance with a unique optimal arm $\as$. \dt{s:Alto}Let
\begin{equation*}
  \Alto(\nu;\mathcal E):=\big\{\nu'\in\Alt(\nu;\mathcal E):\ \nu'_{\as}=\nu_{\as}\big\}
\end{equation*}
be the alternatives that keep the law of the optimal arm. Then
\begin{equation*}
  \GL(\nu;\mathcal E)=\inf\Big\{\sum_{a\ne\as}w_a\gap_a:\ w\in[0,\infty)^{\arms\setminus\{\as\}},\
  \sum_{a\ne\as}w_a\KL(\nu_a,\nu'_a)\ge1\ \text{for every }\nu'\in\Alto(\nu;\mathcal E)\Big\}.
\end{equation*}
\end{lemma}
(see the proof in Appendix~\ref{sec:proof-gl-reduction}). The right-hand side is the formulation
with finite weights and constraints indexed by the alternatives that keep the law of the optimal
arm, as in \citet[Theorem~1]{CMP17}; this constraint set goes back to finite parameter spaces and
controlled Markov chains \citep[see also][]{ATA89,GL97}. Two consequences are immediate: $\GL(\nu;\mathcal E)\le
\GL(\nu;\mathcal E')$ when $\mathcal E\subseteq\mathcal E'$, since the constraints of $\mathcal E$
are among those of $\mathcal E'$; and $\GL(\nu;\mathcal E)=0$ when $\Alto(\nu;\mathcal E)$ is empty,
since $w=0$ is then feasible.

\begin{theorem}[Graves--Lai constant of an unstructured class]\label{thm:main}
\uses{def:gl, def:product, def:kinf}
For every family $(\mathcal D_a)_{a\in\arms}$ of models and every instance $\nu\in\prodc{\mathcal D}$
with a unique optimal arm $\as$,
\begin{equation*}
  \GL\Big(\nu;\prodc{\mathcal D}\Big)=\sum_{a\ne\as}\frac{\gap_a}{\Kinf(\nu_a,\ms;\mathcal D_a)}\in[0,\infty],
\end{equation*}
with $x/0:=\infty$ for $x>0$ and $x/\infty:=0$. In particular, when every model is the same model
$\mathcal D$, $\GL(\nu;\powc{\mathcal D})=\sum_{a\ne\as}\gap_a/\Kinf(\nu_a,\ms;\mathcal D)$.
\end{theorem}
(see the proof in Appendix~\ref{sec:proof-main}).

\paragraph{Interpretation} The Graves--Lai problem lets the exploration rates of all the arms
interact, through alternatives that move several arms at once. On an unstructured class these
alternatives are useless: by Lemma~\ref{lem:gl-reduction} only the alternatives that keep the law
of the optimal arm matter, each of them raises some suboptimal arm $b$ above $\ms$ at a cost of at
least $\Kinf(\nu_b,\ms;\mathcal D_b)$, and the alternatives that move a single arm are already the
binding ones. The optimal rates are $w_a=1/\Kinf(\nu_a,\ms;\mathcal D_a)$, the rates of the
Lai--Robbins bound.

\paragraph{Degenerate values} No regularity of the models is assumed. A suboptimal arm whose
divergence $\Kinf(\nu_a,\ms;\mathcal D_a)$ vanishes makes the constant infinite; by
Lemma~\ref{lem:kinf-p1} this happens on $\powc{\Pone}$ for every instance. An arm whose divergence is
infinite, because no law of $\mathcal D_a$ has a mean above $\ms$, contributes $0$; on $\powc{\Punit}$
this is the case of every arm when $\ms=1$ (Lemma~\ref{lem:kinf-unit}(i)). The proof uses that $\nu$
lies in the class only to build the alternatives that move one arm.

\paragraph{Scope} The theorem concerns unstructured classes and a unique optimal arm. It says
nothing about the attainment of the constant by an algorithm, nor about classes that are not
unstructured, for which the constant is in general not a sum over the arms.

\begin{proposition}[The constant vanishes when the optimal arm identifies the instance]\label{prop:gl-zero}
\uses{def:gl, lem:gl-reduction}
Let $\mathcal E$ be a class and $\nu$ an instance with a unique optimal arm $\as$. If every
$\nu'\in\mathcal E$ with $\nu'_{\as}=\nu_{\as}$ is equal to $\nu$, then $\GL(\nu;\mathcal E)=0$.
\end{proposition}
(see the proof in Appendix~\ref{sec:proof-gl-zero}). For example, on two arms, the class of two
instances
\begin{equation*}
  \mathcal E_2:=\big\{(\Ber(3/4),\Ber(1/4)),\ (\Ber(1/4),\Ber(3/4))\big\},
\end{equation*}
whose optimal arm reveals which instance is played, has the constant $0$ at both instances, while
the unstructured class $\powc{\Bclass}$ has the positive constant $\frac12/\kl(\frac14,\frac34)$ at the
same instances (Theorem~\ref{thm:main} and Lemma~\ref{lem:kinf-unit}(iii)). The class $\mathcal E_2$ is
not unstructured: it does not contain $(\Ber(3/4),\Ber(3/4))$, which any unstructured class containing
it contains. This is the exclusion for
Definition~\ref{def:gl}: the constant is not a function of the instance alone.

The last result of this section connects the constant with the lower bound. We use the form of
the Lai--Robbins and Burnetas--Katehakis bound proved by \citet[Theorem~1]{GMS19}, recalled as
Theorem~\ref{cite:GMS19-thm1} in Appendix~\ref{sec:cited}.

\begin{corollary}[The constant is the lower-bound constant]\label{cor:lower}
\uses{thm:main, def:consistent, def:gl, cite:GMS19-thm1}
Let $\mathcal D$ be a model, $\pi$ a learner consistent on $\powc{\mathcal D}$, and
$\nu\in\powc{\mathcal D}$ an instance with a unique optimal arm. Then, in $[0,\infty]$,
\begin{equation*}
  \liminf_{T\to\infty}\frac{\Reg_T(\pi,\nu)}{\log T}\ge\GL(\nu;\powc{\mathcal D}).
\end{equation*}
\end{corollary}
(see the proof in Appendix~\ref{sec:proof-lower}).

% ----------------------------------------------------------------
\section{Properties of the models and of \texorpdfstring{$\mathcal{K}_{\inf}$}{Kinf}}\label{sec:properties}

In this section we state the basic properties of $\Kinf$: its monotonicity, its
values on laws on $[0,1]$ and on Bernoulli laws, its closed form on Gaussian laws, and its vanishing
on the model of all laws with a finite mean; then the properties of the semi-bounded, moment and
kurtosis models, with a witness and an exclusion for each model. The proofs are in
Appendix~\ref{sec:proof-properties}.

\begin{lemma}[Monotonicity and trivial values]\label{lem:kinf-basic}
\uses{def:kinf}
Let $\mathcal D,\mathcal D'$ be models, $P$ a law and $\mu\le\mu'$ reals. Then:
\begin{enumerate}[label=(\roman*)]
  \item $\Kinf(P,\mu;\mathcal D)\le\Kinf(P,\mu';\mathcal D)$;
  \item if $\mathcal D\subseteq\mathcal D'$, then $\Kinf(P,\mu;\mathcal D')\le\Kinf(P,\mu;\mathcal D)$;
  \item if no law of $\mathcal D$ has a mean larger than $\mu$, then $\Kinf(P,\mu;\mathcal D)=+\infty$;
  \item if $P\in\mathcal D$ and $\mean(P)>\mu$, then $\Kinf(P,\mu;\mathcal D)=0$.
\end{enumerate}
\end{lemma}
(see the proof in Appendix~\ref{sec:proof-kinf-basic}).

\begin{lemma}[$\Kinf$ on laws on {$[0,1]$} and on Bernoulli laws]\label{lem:kinf-unit}
\uses{def:kinf, def:models}
Let $P\in\Punit$ and $\mu\in\R$. Then:
\begin{enumerate}[label=(\roman*)]
  \item if $\mu\ge1$, then $\Kinf(P,\mu;\Punit)=+\infty$;
  \item if $\mean(P)\le\mu<1$, then
  \begin{equation*}
    \kl(\mean(P),\mu)\le\Kinf(P,\mu;\Punit)\le\log\frac{1-\mean(P)}{1-\mu};
  \end{equation*}
  in particular $0<\Kinf(P,\mu;\Punit)<\infty$ when $\mean(P)<\mu<1$, and $\Kinf(P,\mean(P);\Punit)=0$
  when $\mean(P)<1$;
  \item if $P=\Ber(p)$ and $p\le\mu<1$, then
  \begin{equation*}
    \Kinf(\Ber(p),\mu;\Punit)=\Kinf(\Ber(p),\mu;\Bclass)=\kl(p,\mu).
  \end{equation*}
\end{enumerate}
\end{lemma}
(see the proof in Appendix~\ref{sec:proof-kinf-unit}). Item~(iii) is the known case: on Bernoulli
laws, $\Kinf$ is the Bernoulli divergence of the means, and no law on $[0,1]$ does better than a
Bernoulli law. The lower bound of~(ii) is stated by \citet[Section~6.1, after Lemma~1]{CGMMS13} for
the model of all laws on $[0,1]$, through the inequality $\kl(\mean(P),\mean(Q))\le\KL(P,Q)$ on
$\Punit$; we prove it from the data-processing inequality of \citet[Lemma~1]{GMS19}. Item~(ii) is used with $P=\nu_a$ and $\mu=\ms$: on $\powc{\Punit}$, every suboptimal arm
has a finite positive divergence as soon as $\ms<1$. For the upper bound of~(ii), see also the dual
formula of \citet[(5) and Theorem~8]{HT10}: by Jensen's inequality, its objective
$\E_P[\log(1-\lambda(X-\mu))]$ is at most $\log(1+\lambda(\mu-\mean(P)))\le\log\frac{1-\mean(P)}{1-\mu}$
for every $\lambda\in[0,1/(1-\mu)]$.

\begin{lemma}[$\Kinf$ vanishes on all laws with a finite mean]\label{lem:kinf-p1}
\uses{def:kinf}
For every $P\in\Pone$ and every $\mu\in\R$, $\Kinf(P,\mu;\Pone)=0$.
\end{lemma}
(see the proof in Appendix~\ref{sec:proof-kinf-p1}). This is the exclusion for $\Kinf$: without a
constraint on the laws, a small mass placed far away raises the mean at a vanishing cost, and the
lower-bound constant $1/\Kinf$ is infinite. The same phenomenon for sub-Gaussian laws with an
unknown constant is \citet[Exercise~16.4]{LS}.

\begin{lemma}[$\Kinf$ on Gaussian laws]\label{lem:kinf-gauss}
\uses{def:kinf, def:models}
For every $\mu_0\in\R$, $v>0$ and $\mu\ge\mu_0$,
\begin{equation*}
  \Kinf(\Nor(\mu_0,v),\mu;\Gau)=\frac12\log\Big(1+\frac{(\mu-\mu_0)^2}{v}\Big).
\end{equation*}
\end{lemma}
(see the proof in Appendix~\ref{sec:proof-kinf-gauss}). This is the constant
$\frac12\log(1+(\mu^\star-\mu_a)^2/\sigma_a^2)$ of \citet[Section~4.4]{BK96} for normal laws with
unknown mean and variance, restated by \citet[Section~2]{HT14}. In \citet{BK96} the variance $0$ is
allowed in the parameter set, while our model $\Gau$ excludes it.

\begin{lemma}[Semi-bounded and moment models]\label{lem:semi}
\uses{def:kinf, def:models}
Let $P\in\Semi$, $\mu\in\R$, $\varepsilon>0$ and $B>0$. Then:
\begin{enumerate}[label=(\roman*)]
  \item $\Kinf(P,\mu;\Semi)=+\infty$ if and only if $\mu\ge1$; if $\mean(P)\le\mu<1$, then
  $\Kinf(P,\mu;\Semi)\le\log\frac{1-\mean(P)}{1-\mu}$;
  \item for every instance $\nu$ of $\powc{\Semi}$, $\ms<1$ if and only if no arm has the law $\delta_1$;
  \item if $P\in\Punit$, then $\Kinf(P,\mu;\Semi)\le\Kinf(P,\mu;\Punit)$;
  \item every $Q\in\Lmom_{\varepsilon,B}$ has $\int|x|\,dQ(x)\le B^{1/(1+\varepsilon)}$, hence $Q\in\Pone$;
  the law $\delta_0$ belongs to $\Lmom_{\varepsilon,B}$, and $\delta_x$ does not when $|x|^{1+\varepsilon}>B$.
\end{enumerate}
\end{lemma}
(see the proof in Appendix~\ref{sec:proof-semi}). Item~(iii) is only an inequality: the equality,
for a law on $[0,1]$, follows from the dual formula of \citet[Theorem~2]{HT15}, which we do not
prove.

\begin{lemma}[Kurtosis model]\label{lem:kurt}
\uses{def:models, def:kinf}
Let $\kappa>1$. Then:
\begin{enumerate}[label=(\roman*)]
  \item every law $P$ with $\int x^4\,dP<\infty$ and $\var(P)>0$ has $\kurt(P)\ge1$;
  \item let $c\in\R$, $b>0$ and $\phi(x):=c+bx$; every law $P$ with $\int x^4\,dP<\infty$ and
  $\var(P)>0$ has $\kurt(\phi_\#P)=\kurt(P)$; for every law $P$, $\phi_\#P\in\Hk_\kappa$ if and only if
  $P\in\Hk_\kappa$; and
  \begin{equation*}
    \Kinf(\phi_\#P,\phi(\mu);\Hk_\kappa)=\Kinf(P,\mu;\Hk_\kappa)\qquad\text{for every law }P\text{ and every }\mu\in\R;
  \end{equation*}
  \item (witnesses) $\Ber(1/2)\in\Hk_\kappa$, and $\Nor(\mu_0,v)\in\Hk_\kappa$ if and only if $\kappa\ge3$;
  \item (exclusions) $\delta_x\notin\Hk_\kappa$ for every $x$, and $\Ber(p)\notin\Hk_\kappa$ when
  $0<p<1$ and $p(1-p)<1/(\kappa+3)$.
\end{enumerate}
\end{lemma}
(see the proof in Appendix~\ref{sec:proof-kurt}). Item~(ii) is the scale invariance that makes the
kurtosis model the natural model without a known scale \citep[Assumption~1]{L17}.

% ----------------------------------------------------------------
\section{Classes of learners: relations, witnesses and exclusions}\label{sec:learners}

In this section we relate the classes of learners, show that the rewards collected by count form an
independent array in every run, give for each class a learner in it and a learner outside it, and
define the further classes of the lower bounds in the initial regime and of the tails.
The proofs are in Appendix~\ref{sec:proof-learners}.

\begin{lemma}[Relations between the classes]\label{lem:classes}
\uses{def:consistent, def:ufc, def:smarter}
Let $\mathcal E\subseteq\mathcal E'$ be classes, $\pi$ a learner, $C>0$ and $0\le\gamma\le\gamma'<1$.
Then:
\begin{enumerate}[label=(\roman*)]
  \item each property of Definitions~\ref{def:consistent}--\ref{def:smarter} on $\mathcal E'$
  implies the same property on $\mathcal E$;
  \item $\gamma$-consistency on $\mathcal E$ implies $\gamma'$-consistency on $\mathcal E$;
  \item consistency on $\mathcal E$ implies uniform fast convergence on $\mathcal E$;
  \item uniform super-fast convergence on $\mathcal E$ with the constant $C$ implies consistency on
  $\mathcal E$ and uniform fast convergence on $\mathcal E$;
  \item if $\pi$ is consistent on $\mathcal E$, the constant horizon-dependent learner
  $(\pi)_{T\ge1}$ is consistent on $\mathcal E$.
\end{enumerate}
\end{lemma}
(see the proof in Appendix~\ref{sec:proof-classes}).

The following lemma is the bridge between runs, defined by conditional laws, and the ``stack of
rewards'' in which regret proofs are usually written: the $s$-th reward collected from arm $a$ has
the law $\nu_a$, independently of everything else collected.

\begin{lemma}[Rewards by count]\label{lem:stack}
\uses{def:run, def:regret}
Consider a run of a learner $\pi$ on an instance $\nu$, on $(\Omega,\mathcal F,\Prob)$. Enlarge it
to $\Omega\times\R^{\arms\times\mathbb N_{\ge1}}$ with the law $\Prob\otimes\bigotimes_{(a,s)}\nu_a$,
\dt{s:Z}whose coordinates $\Zc_{a,s}$ are independent of the run, with $\Zc_{a,s}\sim\nu_a$.
\dt{s:tau}For $a\in\arms$ and $s\ge1$, let $\taup_{a,s}:=\inf\{t\ge0:\Np_a(t+1)=s\}\in\mathbb N\cup\{\infty\}$
be the round of the $s$-th pull of $a$, and \dt{s:Y}let
\begin{equation*}
  \Yc_{a,s}:=\Xr_{\taup_{a,s}}\ \text{ if }\taup_{a,s}<\infty,\qquad\Yc_{a,s}:=\Zc_{a,s}\ \text{ otherwise}.
\end{equation*}
Then the variables $(\Yc_{a,s})_{a\in\arms,\,s\ge1}$ are independent and $\Yc_{a,s}$ has the law $\nu_a$.
Moreover, \dt{s:bY}with $\bY_{a,s}:=\frac1s\sum_{j=1}^s\Yc_{a,j}$, for every $t\ge0$ and every arm $a$
with $\Np_a(t)\ge1$,
\begin{equation*}
  \hm_a(t)=\bY_{a,\Np_a(t)}.
\end{equation*}
\end{lemma}
(see the proof in Appendix~\ref{sec:proof-stack}). The enlargement changes neither the law of the
run nor any expectation $\Er_{\pi,\nu}$. The lemma restates a known fact, a form of the optional
skipping theorem \citep[see also][]{D36}; our proof by truncation in time is the one of its
formalisation in \citet{LML} (\texttt{RewardByCountMeasure.lean}), and \citet[Section~4.6]{LS}
assert the equivalence of the two models.

We now define the witness of the convergence classes: the upper confidence bound algorithm with
the bonus given by Hoeffding's inequality. Its index is, at round $t$,
\begin{equation}\label{eq:ucb}
  \dt{s:U}\Uix_a(t):=\hm_a(t)+\sqrt{\frac{2\log(t+1)}{\Np_a(t)}}\quad\text{if }\Np_a(t)\ge1,\qquad
  \Uix_a(t):=+\infty\quad\text{if }\Np_a(t)=0,
\end{equation}
a function of the history $\Hist_t$: an arm is optimistic about its mean by the half-width of a
Hoeffding confidence interval at level $(t+1)^{-4}$.

\begin{algorithm}[t]
\caption{\texttt{UCB} (Hoeffding bonus, any tie rule)}\label{alg:ucb}
\begin{algorithmic}[1]
\Require the arms $\arms$; a tie rule $(\rho_t)_{t\ge0}$
\For{$t = 0, 1, 2, \dots$}
  \State compute $\Uix_a(t)$ for every arm $a$ by \eqref{eq:ucb}
  \State draw $\Ap_t$ from $\rho_t(\Hist_t,\cdot)$, a law on $\argmax_{a\in\arms}\Uix_a(t)$
  \State observe $\Xr_t$
\EndFor
\end{algorithmic}
\end{algorithm}

\begin{definition}[\texttt{UCB}]\label{def:ucb}
\uses{def:learner, def:regret}
A \emph{tie rule} is a sequence of Markov kernels $\rho_t$ from $\Hs_t$ to $\arms$ such that
$\rho_t(h,M_t(h))=1$ for every $h\in\Hs_t$, where $M_t(h):=\argmax_{a\in\arms}\Uix_a(t)$ is computed
by \eqref{eq:ucb} from the history $h$. The learner \ucb{} with the tie rule $(\rho_t)$ is
$\pi=(\rho_t)_{t\ge0}$ (Algorithm~\ref{alg:ucb}).
\end{definition}

Every arm is pulled once before any arm is pulled twice, since an arm never pulled has the index
$+\infty$. The uniform law on $M_t(h)$ is a tie rule; the learner is anytime.

\begin{proposition}[\texttt{UCB} is uniformly super-fast convergent on laws on {$[0,1]$}]\label{prop:ucb}
\uses{def:ucb, def:ufc, def:consistent, def:models}
For every tie rule, every instance $\nu\in\powc{\Punit}$, every suboptimal arm $a$ and every $T\ge1$,
the learner \ucb{} satisfies
\begin{equation*}
  \Er[\Np_a(T)]\le\frac{8\log T}{\gap_a^2}+\frac{2}{\gap_a^2}+\frac54,
\end{equation*}
and, for every $T\ge2$, $\Er[\Np_a(T)]\le13\log T/\gap_a^2$. Hence \ucb{} is uniformly super-fast
convergent on $\powc{\Punit}$ with the constant $13$, consistent on $\powc{\Punit}$ and on
$\powc{\Bclass}$, and uniformly fast convergent on both.
\end{proposition}
(see the proof in Appendix~\ref{sec:proof-ucb}). This is the analysis of UCB1 by
\citet[Theorem~1]{ACF02}, whose index is the empirical mean plus $\sqrt{2\ln n/n_j}$, with $n$ the
number of plays so far, and whose proof gives $\E[T_j(n)]\le8\ln n/\gap_j^2+1+\pi^2/3$ for each
suboptimal arm. The leading term $8\log T/\gap_a^2$ is the same; our second term
$2/\gap_a^2+5/4$, which comes from a sum over the samples of $a$ alone, is larger than theirs when
$\gap_a$ is small. The point here is that the super-fast class is not empty on $[0,1]$-valued
instances, with a constant computed in our conventions. KL-UCB is consistent
too, on Bernoulli instances with means in $(0,1)$, by the finite-time bound of
\citet[Theorem~1]{CGMMS13}; we do not use it.

\begin{lemma}[Constant and uniform learners]\label{lem:constant}
\uses{def:consistent, def:ufc, def:smarter, def:regret, def:models}
\dt{s:pib}For $b\in\arms$, let $\Cons^b$ be the learner that always pulls $b$ ($\pol_t(h,\cdot)=\delta_b$),
and let $\pi^{\mathrm{u}}$ be the uniform learner ($\pol_t(h,\cdot)$ uniform on $\arms$). For every
instance $\nu$ and every $T\ge1$:
\begin{enumerate}[label=(\roman*)]
  \item $\Er_{\Cons^b,\nu}[\Np_b(T)]=T$ and $\Reg_T(\Cons^b,\nu)=\gap_bT$;
  \item $\Er_{\pi^{\mathrm{u}},\nu}[\Np_a(T)]=T/\Karms$ for every arm $a$, and $\Reg_T(\pi^{\mathrm{u}},\nu)=\frac{T}{\Karms}\sum_a\gap_a$.
\end{enumerate}
Consequently, on every class $\mathcal E$ that contains an instance in which $b$ is suboptimal (for
example $\powc{\Bclass}$), $\Cons^b$ is neither $\gamma$-consistent for any $\gamma\in[0,1)$, nor uniformly
fast convergent, nor uniformly super-fast convergent with any constant, nor smarter than the
uniform learner; and $\pi^{\mathrm{u}}$ is smarter than the uniform learner on every class, and not
$\gamma$-consistent for any $\gamma\in[0,1)$ on a class that contains an instance with a suboptimal arm.
\end{lemma}
(see the proof in Appendix~\ref{sec:proof-constant}).

\subsection{Further classes of learners}

The lower bounds in the initial regime and the results on the distribution of the regret use
further classes, which the main result does not use. We define them here, on an arbitrary class, with
a witness and an exclusion for each. The first two are from \citet[Definitions~3 and~4]{GMS19}.

\begin{definition}[Pairwise symmetric for optimal arms]\label{def:pairwise}
\uses{def:regret, def:product}
Let $\mathcal E$ be a class. \dt{w:pairwise}A learner $\pi$ is \emph{pairwise symmetric for optimal
arms on $\mathcal E$} if, for every $\nu\in\mathcal E$, every pair of optimal arms $a,a'\in\Opt(\nu)$
with $\nu_a=\nu_{a'}$ and every $T\ge1$, the pairs $(\Np_a(T),\Np_{a'}(T))$ and
$(\Np_{a'}(T),\Np_a(T))$ have the same law under $\Pru_{\pi,\nu}$.
\end{definition}

\dt{s:preceq}For two instances $\nu,\nu'$, write $\nu'\preccurlyeq\nu$ if $\nu'_a=\nu_a$ for every
$a\in\Opt(\nu)$ and $\marm_a(\nu')\le\marm_a(\nu)$ for every $a\notin\Opt(\nu)$: the optimal arms keep
their laws and the other arms get worse means. Then $\ms(\nu')=\ms(\nu)$ and
$\Opt(\nu')=\Opt(\nu)$, since an arm $a\notin\Opt(\nu)$ has $\marm_a(\nu')\le\marm_a(\nu)<\ms(\nu)$.

\begin{definition}[Monotonic learners]\label{def:monotonic}
\uses{def:regret, def:product}
Let $\mathcal E$ be a class. \dt{w:monotonic}A learner $\pi$ is \emph{monotonic on $\mathcal E$} if,
for all $\nu,\nu'\in\mathcal E$ with $\nu'\preccurlyeq\nu$ and every $T\ge1$,
\begin{equation*}
  \sum_{a\in\Opt(\nu')}\Er_{\pi,\nu'}[\Np_a(T)]\ge\sum_{a\in\Opt(\nu)}\Er_{\pi,\nu}[\Np_a(T)].
\end{equation*}
\end{definition}

A monotonic learner pulls the optimal arms at least as often, in total, when the other arms get
worse. \citet[Definition~4]{GMS19} write the inequality without a quantifier on $T$; we require it
for every $T\ge1$.

The tail classes are from \citet[Definition~1]{SA11}. \dt{s:dmin}For an instance $\nu$ with a unique
optimal arm $\as$, let $\gap_{\min}:=\min_{a\ne\as}\gap_a>0$ be the smallest gap.

\begin{definition}[Upper-tailed counts and pseudo-regret]\label{def:tails}
\uses{def:regret, def:product}
Let $\mathcal E$ be a class and $f:\{2,3,\dots\}\to(0,\infty)$. \dt{w:tailT}A learner $\pi$ has
\emph{$f$-upper tailed counts on $\mathcal E$} if there are $C,\tilde C>0$ such that, for every
$\nu\in\mathcal E$ with a unique optimal arm, every suboptimal arm $a$ and every $n\ge2$,
\begin{equation*}
  \Pru_{\pi,\nu}\Big(\Np_a(n)\ge\frac{C\log n}{\gap_a^2}\Big)\le\frac{\tilde C}{f(n)}.
\end{equation*}
\dt{w:tailR}It has \emph{$f$-upper tailed pseudo-regret on $\mathcal E$} if there are $C,\tilde C>0$
such that, for every $\nu\in\mathcal E$ with a unique optimal arm and every $n\ge2$,
\begin{equation*}
  \Pru_{\pi,\nu}\Big(\hR_n\ge\frac{C\log n}{\gap_{\min}}\Big)\le\frac{\tilde C}{f(n)}.
\end{equation*}
\end{definition}

\paragraph{Changes to the source} \citet[Definition~1]{SA11} take a class of instances whose laws are
on $[0,1]$, and the regret against the rewards of the optimal arm,
$\sum_{t=1}^nX_{\as,t}-\sum_{t=1}^nX_{I_t,T_{I_t}(t)}$ in their notation, which is not a function of
the counts. We take an arbitrary class and the pseudo-regret $\hR_n=\sum_a\gap_a\Np_a(n)$ of
Definition~\ref{def:regret}, so that Definition~\ref{def:tails} for the pseudo-regret is a
property of the counts (Lemma~\ref{lem:tails}(i)). Their counts are ours: $T_k(n)$, the number of
pulls of $k$ among the rounds $1,\dots,n$, is our $\Np_k(n)$. Their weak versions
\citep[Definition~3]{SA11}, in which $C$ and $\tilde C$ may depend on the instance, are not used here.

\begin{definition}[Horizon-dependent versions]\label{def:family}
\uses{def:consistent, def:ufc, def:smarter, def:pairwise, def:monotonic, def:tails}
Let $\mathcal E$ be a class and $(\pi^T)_{T\ge1}$ a horizon-dependent learner. \dt{w:family}The
family is \emph{$\gamma$-consistent}, \emph{uniformly fast convergent}, \emph{uniformly super-fast
convergent with the constant $C$}, \emph{smarter than the uniform learner}, \emph{pairwise symmetric
for optimal arms}, \emph{monotonic}, or has \emph{$f$-upper tailed counts} or \emph{pseudo-regret} on
$\mathcal E$, if the condition of Definition~\ref{def:consistent}, \ref{def:ufc}, \ref{def:smarter},
\ref{def:pairwise}, \ref{def:monotonic} or~\ref{def:tails} holds when, at each horizon $T$ (at each
$n$ in Definition~\ref{def:tails}), the law $\Pru_{\pi,\nu}$ and the expectation $\Er_{\pi,\nu}$ of the
counts and of the regret at that horizon are those of $\pi^T$ (respectively $\pi^n$).
\end{definition}

For consistency this is the definition of a consistent horizon-dependent learner in
Definition~\ref{def:consistent}. The witnesses and exclusions are the following.

\begin{lemma}[Pairwise symmetric and monotonic learners]\label{lem:gms-classes}
\uses{def:pairwise, def:monotonic, def:models, def:learner, lem:constant}
Let $b\ne c$ be two arms, $\Cons^b$ and $\pi^{\mathrm{u}}$ the learners of Lemma~\ref{lem:constant},
\dt{s:pisw}and $\pi^{\mathrm{sw}}$ the learner that pulls $b$ at round $0$ and then, at every round
$t\ge1$, pulls $b$ if $\Xr_0\le1/2$ and $c$ if $\Xr_0>1/2$. Then:
\begin{enumerate}[label=(\roman*)]
  \item $\pi^{\mathrm{u}}$ is pairwise symmetric for optimal arms and monotonic on every class;
  \item $\Cons^b$ is not pairwise symmetric for optimal arms on any class that contains an instance
  $\nu$ in which $b$ and another arm $a$ are optimal with $\nu_a=\nu_b$ (for example $\powc{\Bclass}$);
  \item $\pi^{\mathrm{sw}}$ is not monotonic on $\powc{\Bclass}$.
\end{enumerate}
\end{lemma}
(see the proof in Appendix~\ref{sec:proof-gms-classes}).

\begin{lemma}[Tails of the counts and of the pseudo-regret]\label{lem:tails}
\uses{def:tails, def:regret, def:models, lem:constant}
Let $\mathcal E$ be a class, $f:\{2,3,\dots\}\to(0,\infty)$ and $\pi$ a learner. Then:
\begin{enumerate}[label=(\roman*)]
  \item if $\pi$ has $f$-upper tailed counts on $\mathcal E$ with the constants $C,\tilde C$, it has
  $f$-upper tailed pseudo-regret on $\mathcal E$ with the constants $(\Karms-1)C$ and $(\Karms-1)\tilde C$;
  \item if $f$ is unbounded and $\mathcal E$ contains an instance with a unique optimal arm in which
  $b$ is suboptimal (for example $\powc{\Bclass}$), then $\Cons^b$ has neither $f$-upper tailed counts
  nor $f$-upper tailed pseudo-regret on $\mathcal E$.
\end{enumerate}
\end{lemma}
(see the proof in Appendix~\ref{sec:proof-tails}).

\begin{proposition}[Tails of the counts of \texttt{UCB}]\label{prop:ucb-tails}
\uses{def:ucb, def:tails, def:models, lem:stack, lem:tails}
For every tie rule, every instance $\nu\in\powc{\Punit}$, every suboptimal arm $a$ and every
$T\ge2$, the learner \ucb{} satisfies
\begin{equation*}
  \Pru\Big(\Np_a(T)\ge\frac{10\log T}{\gap_a^2}\Big)\le\frac{1}{(\log T)^2}.
\end{equation*}
Hence \ucb{} has $f$-upper tailed counts on $\powc{\Punit}$ for $f(n):=(\log n)^2$, with the
constants $10$ and $1$, and $f$-upper tailed pseudo-regret with the constants $10(\Karms-1)$ and
$\Karms-1$.
\end{proposition}
(see the proof in Appendix~\ref{sec:proof-ucb-tails}). We do not optimise the tail
$1/(\log T)^2$: \citet{SA11} study how fast the tail of an anytime learner can decrease.

\begin{lemma}[Constant families]\label{lem:family}
\uses{def:family}
Let $\mathcal E$ be a class and $\pi$ a learner. The constant horizon-dependent learner
$(\pi)_{T\ge1}$ has a property of Definition~\ref{def:family} on $\mathcal E$ if and only if $\pi$ has
the same property on $\mathcal E$.
\end{lemma}
(see the proof in Appendix~\ref{sec:proof-family}). With Proposition~\ref{prop:ucb} and
Lemma~\ref{lem:constant}, the constant family of \ucb{} is uniformly super-fast convergent on
$\powc{\Punit}$ with the constant $13$, and the constant family of $\Cons^b$ is not uniformly fast
convergent on $\powc{\Bclass}$.

% ----------------------------------------------------------------
\section{How the proofs work}\label{sec:overview}

Most proofs are short; we describe the three that are not, and say where the degenerate values are
handled.

\paragraph{The identity on unstructured classes} The proof of Theorem~\ref{thm:main} has two
steps. Lemma~\ref{lem:gl-reduction} removes the optimal arm: a feasible weight vector of the
reduced problem becomes feasible for the full problem by setting $w_{\as}:=\infty$, at no cost
since $\gap_{\as}=0$, because an alternative that changes the law of the optimal arm has
$\KL(\nu_{\as},\nu'_{\as})>0$ and its constraint is then infinite. Conversely, a feasible vector of
the full problem restricts to a feasible vector of the reduced one, since the term of $\as$ vanishes
on the alternatives that keep $\nu_{\as}$. Then, on $\prodc{\mathcal D}$, the lower bound uses the
alternatives that replace one suboptimal arm $a$ by a law $Q\in\mathcal D_a$ with $\mean(Q)>\ms$:
their constraints read $w_a\KL(\nu_a,Q)\ge1$, which force $w_a\ge1/\Kinf(\nu_a,\ms;\mathcal D_a)$. For
the upper bound, the weights $1/\Kinf$ are feasible: an alternative that keeps $\nu_{\as}$ has some
arm $b$ with mean above $\ms$, so its constraint holds through the single term
$\KL(\nu_b,\nu'_b)\ge\Kinf(\nu_b,\ms;\mathcal D_b)$. The degenerate values $\Kinf\in\{0,\infty\}$ are
checked case by case, with the conventions of $[0,\infty]$.

\paragraph{Rewards by count} Lemma~\ref{lem:stack} is proved by truncation in time. The array
$\Yc^t$ in which the rewards collected at rounds $u\ge t$ are replaced by the independent
variables $\Zc$ has the product law at $t=0$; passing from $\Yc^t$ to $\Yc^{t+1}$ replaces a single
coordinate, the one of the pull of round $t$, whose conditional law given the rest is $\nu_{\Ap_t}$
in both arrays. So all the $\Yc^t$ have the product law, and $\Yc^t\to\Yc$ coordinatewise.

\paragraph{The analysis of \texttt{UCB}} Proposition~\ref{prop:ucb} follows the classical
decomposition: the pulls of a suboptimal arm $a$ before it has $u\approx8\log T/\gap_a^2$ samples,
the rounds at which an optimal arm is underestimated, and the samples $s\ge u$ at which the
empirical mean of $a$ exceeds $\marm_a+\gap_a/2$. Lemma~\ref{lem:ucb-decomp} states the
decomposition and Lemma~\ref{lem:ucb-terms} bounds the two random terms by Hoeffding's inequality,
applied to the array of Lemma~\ref{lem:stack}. Proposition~\ref{prop:ucb-tails} uses the same
argument at a single count $s\approx10\log T/\gap_a^2$: if $a$ is pulled at count $s$, either the
empirical mean of its first $s$ samples is too large, or an optimal arm is underestimated at a round
after $s$, and both events have small probability.

\paragraph{Map of the appendices} Appendix~\ref{sec:cited} recalls the two results we cite.
Appendix~\ref{sec:proof-gl} proves the results on the Graves--Lai constant,
Appendix~\ref{sec:proof-properties} the properties of $\Kinf$ and of the models, and
Appendix~\ref{sec:proof-learners} those of the classes of learners. Appendix~\ref{sec:attribution}
credits the techniques and Appendix~\ref{sec:notation} lists the symbols.

% ----------------------------------------------------------------
\section{Examples}\label{sec:examples}

In this section we compute the constants on explicit instances. The numbers can be reproduced with
a short script distributed with the source of this paper. No simulation is reported: the only
algorithm of the paper, \ucb{}, is a witness of a class, and its bound is far from the constant of
the lower bound (on the Bernoulli instance below, $13/\gap_a^2=325$ against
$1/\kl(0.5,0.7)\approx11.5$), which a simulation would only confirm.

\begin{example}[Bernoulli instance]
Let $\arms=\{1,2,3\}$ and $\nu=(\Ber(0.7),\Ber(0.5),\Ber(0.3))\in\powc{\Bclass}$. Its optimal arm is
$1$, with $\ms=0.7$, $\gap_2=0.2$ and $\gap_3=0.4$. By Lemma~\ref{lem:kinf-unit}(iii) and
Theorem~\ref{thm:main},
\begin{equation*}
  \GL(\nu;\powc{\Bclass})=\GL(\nu;\powc{\Punit})=\frac{0.2}{\kl(0.5,0.7)}+\frac{0.4}{\kl(0.3,0.7)}
  \approx\frac{0.2}{0.0872}+\frac{0.4}{0.3389}\approx3.474.
\end{equation*}
The constants of the two classes coincide, although $\powc{\Punit}$ is much larger: the Bernoulli laws
are the hardest laws on $[0,1]$ with given means.
\end{example}

\begin{example}[Gaussian instance, unknown variance]
Let $\nu=(\Nor(0,1),\Nor(-1,1))\in\powc{\Gau}$. By Lemma~\ref{lem:kinf-gauss},
$\Kinf(\Nor(-1,1),0;\Gau)=\frac12\log2\approx0.3466$, so $\GL(\nu;\powc{\Gau})=2/\log2\approx2.885$.
With a known variance, the model $\{\Nor(\mu,1):\mu\in\R\}$ gives
$\Kinf=\frac12$ and the constant $2$: not knowing the variance costs a factor $1/\log2\approx1.44$ on
this instance.
\end{example}

\begin{example}[A law on {$[0,1]$} that is not Bernoulli]
Let $P:=0.2\,\delta_0+0.5\,\delta_{1/2}+0.3\,\delta_1$, with $\mean(P)=0.55$, and $\mu=0.8$. By
Lemma~\ref{lem:kinf-unit}(ii),
\begin{equation*}
  0.1588\approx\kl(0.55,0.8)\le\Kinf(P,0.8;\Punit)\le\log\frac{0.45}{0.2}\approx0.8109.
\end{equation*}
\end{example}

\paragraph{Degenerate cases} Let $\nu\in\powc{\Punit}$ have a unique optimal arm. If $\ms=1$, every
divergence $\Kinf(\nu_a,1;\Punit)$ is infinite and $\GL(\nu;\powc{\Punit})=0$
(Lemma~\ref{lem:kinf-unit}(i) and Theorem~\ref{thm:main}). On $\powc{\Pone}$ the
constant is infinite at every instance with a unique optimal arm (Lemma~\ref{lem:kinf-p1}), so by
Corollary~\ref{cor:lower} every learner consistent on $\powc{\Pone}$, if there is one, has a regret
that grows faster than any multiple of $\log T$. With $\Karms=2$ and $\nu=(\Ber(p),\Ber(q))$, $p>q$,
the constant is $(p-q)/\kl(q,p)$ for $p<1$, and $0$ for $p=1$.

% ----------------------------------------------------------------
\section{Conclusion and open questions}\label{sec:open}

We fixed the definitions of the stochastic bandit setting with finitely many arms, proved their
basic properties with every degenerate value, gave a witness and an exclusion for each, and proved
that the Graves--Lai constant of an unstructured class is the sum of the Lai--Robbins rates.

\paragraph{Limitations} Some objects of the same setting are not defined here: the algorithms
analysed by other papers, which are defined where they are analysed; the dual formulas of $\Kinf$
on $[0,1]$ and on semi-bounded laws \citep[Theorem~8]{HT10}, \citep[Theorem~2]{HT15}; the weak tail
classes of \citet[Definition~3]{SA11}; and a definition of the Graves--Lai constant with several
optimal arms. The tail of \ucb{} in Proposition~\ref{prop:ucb-tails} is only a witness that the tail
classes are not empty for a growing $f$.

\paragraph{Open questions} None: this paper fixes definitions.

% ================================================================ appendices
\appendix

\section{Useful lemmas}\label{sec:cited}

In this section we present the two results from other papers that we use: the data-processing
inequality for a $[0,1]$-valued variable, used for the lower bound of Lemma~\ref{lem:kinf-unit}(ii),
and the lower bound of \citet[Theorem~1]{GMS19}, used in Corollary~\ref{cor:lower}.

\citet{GMS19} write $\mathbb P_1,\mathbb P_2$ for two distributions on a measurable space
$(\Gamma,\mathcal G)$ and $\mathbb E_1,\mathbb E_2$ for their expectations; the other symbols are
ours.

\begin{lemma}[{From \citet[Lemma~1]{GMS19}}]\label{cite:GMS19-lem1}
Let $(\Gamma,\mathcal G)$ be a measurable space with two probability measures $\mathbb P_1$ and
$\mathbb P_2$, and $Z:\Gamma\to[0,1]$ a $\mathcal G$-measurable function. Then
\begin{equation*}
  \KL(\mathbb P_1,\mathbb P_2)\ge\kl\big(\mathbb E_1[Z],\mathbb E_2[Z]\big).
\end{equation*}
\end{lemma}

\citet{GMS19} write $\underline\nu$ for an instance, $\psi$ for a strategy, $N_{\psi,a}(T)$ for the
number of pulls of $a$ among the rounds $1,\dots,T$ (our $\Np_a(T)$), $\mu^\star$ for our $\ms$,
$E(\cdot)$ for our $\mean(\cdot)$ and $\mathcal K_{\inf}(\nu_a,x,\mathcal D)$ for our
$\Kinf(\nu_a,x;\mathcal D)$. Their strategies draw the arm of each round as a measurable function of
the past rewards and of independent uniform variables; every learner of Definition~\ref{def:learner}
is realised in this way with the same law of $(\Ap_t,\Xr_t)_{t\ge0}$ (draw $\Ap_t$ from
$\pol_t(\Hist_t,\cdot)$ by the quantile transform of a fresh uniform variable, $\arms$ being finite),
so their counts have our expectations. Their ``uniformly fast convergent on $\mathcal D$'' is our
uniform fast convergence on $\powc{\mathcal D}$ (Definition~\ref{def:ufc}). They state the result for
all bandit problems; we use it only for the instances of $\powc{\mathcal D}$.

\begin{theorem}[{From \citet[Theorem~1]{GMS19}}]\label{cite:GMS19-thm1}
\uses{def:ufc, def:kinf}
For every model $\mathcal D$, every learner $\pi$ uniformly fast convergent on $\powc{\mathcal D}$, every
instance $\nu\in\powc{\mathcal D}$ and every suboptimal arm $a$ of $\nu$,
\begin{equation*}
  \liminf_{T\to\infty}\frac{\Er_{\pi,\nu}[\Np_a(T)]}{\log T}\ge\frac{1}{\Kinf(\nu_a,\ms;\mathcal D)}.
\end{equation*}
\end{theorem}

We also use, by name, the following textbook facts: Hoeffding's inequality for independent
variables with values in $[0,1]$ \citep{H63}, the formula for the Kullback--Leibler divergence
between two Gaussian laws, its invariance under measurable bijections, Gibbs' inequality
($\KL(P,Q)=0$ only if $P=Q$), and the Ionescu-Tulcea theorem.

% ----------------------------------------------------------------
\section{Proofs for the Graves--Lai constant}\label{sec:proof-gl}

In this appendix we prove the results of Section~\ref{sec:main}: the reduction, the main theorem,
the vanishing of the constant, and the lower bound. All computations are in $[0,\infty]$ with
$0\cdot\infty=0$. We use repeatedly that, for $w\in[0,\infty]$ and $k\in[0,\infty]$,
\begin{equation}\label{eq:prod}
  w\cdot k=\infty\quad\text{whenever }w=\infty\text{ and }k>0,\qquad w\cdot0=0\quad\text{for every }w.
\end{equation}

\subsection{Proof of Lemma~\ref{lem:gl-reduction}}\label{sec:proof-gl-reduction}

\begin{proof}
\uses{def:gl}
\emph{Setting.} Let $V$ be the right-hand side of the lemma, and call \emph{reduced-feasible} a
vector $w\in[0,\infty)^{\arms\setminus\{\as\}}$ that satisfies its constraints. \emph{Goal.}
$\GL(\nu;\mathcal E)=V$.

\emph{Step 1: $\GL(\nu;\mathcal E)\le V$.} Let $w$ be reduced-feasible and extend it by
$w_{\as}:=\infty$. Let $\nu'\in\Alt(\nu;\mathcal E)$. If $\nu'_{\as}\ne\nu_{\as}$, then
$\KL(\nu_{\as},\nu'_{\as})>0$ by Gibbs' inequality, so the term of $\as$ is $\infty$ by
\eqref{eq:prod} and the constraint holds. If $\nu'_{\as}=\nu_{\as}$, then $\nu'\in\Alto(\nu;\mathcal E)$
and the constraint holds by reduced feasibility, the term of $\as$ being nonnegative. So the
extension is feasible, with the cost
\begin{equation*}
  \sum_{a\in\arms}w_a\gap_a=\infty\cdot0+\sum_{a\ne\as}w_a\gap_a=\sum_{a\ne\as}w_a\gap_a.
\end{equation*}
Taking the infimum over reduced-feasible $w$ gives $\GL(\nu;\mathcal E)\le V$ (if there is no
reduced-feasible $w$, then $V=\infty$).

\emph{Step 2: $V\le\GL(\nu;\mathcal E)$.} Let $w\in[0,\infty]^{\arms}$ be feasible. If $w_a=\infty$
for some $a\ne\as$, its cost is $\infty\cdot\gap_a=\infty$, since $\gap_a>0$ (the optimal arm is
unique), and there is nothing to prove. Otherwise the restriction
$\tilde w:=(w_a)_{a\ne\as}\in[0,\infty)^{\arms\setminus\{\as\}}$ is reduced-feasible: for
$\nu'\in\Alto(\nu;\mathcal E)$, $\KL(\nu_{\as},\nu'_{\as})=\KL(\nu_{\as},\nu_{\as})=0$, so by
\eqref{eq:prod}
\begin{equation*}
  \sum_{a\ne\as}\tilde w_a\KL(\nu_a,\nu'_a)=\sum_{a\in\arms}w_a\KL(\nu_a,\nu'_a)\ge1.
\end{equation*}
Its cost is $\sum_{a\ne\as}w_a\gap_a=\sum_aw_a\gap_a$, since $w_{\as}\gap_{\as}=0$. Hence
$V\le\sum_aw_a\gap_a$, and taking the infimum over feasible $w$ gives $V\le\GL(\nu;\mathcal E)$.
\end{proof}

\subsection{Proof of Theorem~\ref{thm:main}}\label{sec:proof-main}

\begin{proof}
\uses{def:gl, def:product, def:kinf, lem:gl-reduction}
\emph{Idea.} By Lemma~\ref{lem:gl-reduction} (\emph{removing the optimal arm}) we may work with
finite weights on the suboptimal arms and the alternatives that keep $\nu_{\as}$. The alternatives
that move one suboptimal arm give the lower bound; the weights $1/\Kinf$ are feasible, which gives
the upper bound.

\emph{Setting.} $\mathcal E:=\prodc{\mathcal D}$, $\nu\in\mathcal E$ with the unique optimal arm
$\as$, so $\gap_a>0$ for $a\ne\as$. \dt{s:Kd}For $a\ne\as$, write $\Kd_a:=\Kinf(\nu_a,\ms;\mathcal D_a)\in[0,\infty]$, and
$S:=\sum_{a\ne\as}\gap_a/\Kd_a$. Let $V$ be the right-hand side of Lemma~\ref{lem:gl-reduction}, so
that $\GL(\nu;\mathcal E)=V$. \emph{Goal.} $V=S$.

\emph{Step 1: every reduced-feasible $w$ has $w_a\ge1/\Kd_a$ for every $a\ne\as$.} Fix $a\ne\as$ and
$Q\in\mathcal D_a$ with $\mean(Q)>\ms$, and let $\nu'$ be $\nu$ with the law of arm $a$ replaced by $Q$.
Then $\nu'\in\mathcal E$, because $\nu_b\in\mathcal D_b$ for $b\ne a$ (here we use $\nu\in\mathcal E$)
and $Q\in\mathcal D_a$; and $\marm_a(\nu')=\mean(Q)>\ms=\marm_{\as}(\nu')$, so $\as\notin\Opt(\nu')$
and $\nu'_{\as}=\nu_{\as}$, that is $\nu'\in\Alto(\nu;\mathcal E)$. The constraint of $\nu'$ has a
single nonzero term, since $\KL(\nu_b,\nu'_b)=0$ for $b\ne a$:
\begin{equation*}
  w_a\KL(\nu_a,Q)\ge1.
\end{equation*}
As $w_a<\infty$ and $0\cdot\KL(\nu_a,Q)=0$, this forces $w_a>0$ and $\KL(\nu_a,Q)\ge1/w_a$. Taking
the infimum over such $Q$ gives $\Kd_a\ge1/w_a>0$, that is $w_a\ge1/\Kd_a$. If there is no such $Q$,
then $\Kd_a=\infty$ and $w_a\ge0=1/\Kd_a$ holds trivially. So in every case
\begin{equation*}
  w_a\ge\frac{1}{\Kd_a}\qquad\text{(with }1/0:=\infty\text{ and }1/\infty:=0\text{)}.
\end{equation*}

\emph{Step 2: $V\ge S$.} Let $w$ be reduced-feasible. By Step 1, $w_a\ge1/\Kd_a$ for every $a\ne\as$.
If some $\Kd_a=0$, no finite $w_a$ satisfies this, so there is
no reduced-feasible $w$, and $V=\infty=S$ (the term $\gap_a/0=\infty$ makes $S$ infinite). Otherwise
\begin{equation*}
  \sum_{a\ne\as}w_a\gap_a\ge\sum_{a\ne\as}\frac{\gap_a}{\Kd_a}=S,
\end{equation*}
and taking the infimum over $w$ gives $V\ge S$.

\emph{Step 3: $V\le S$.} If $S=\infty$ there is nothing to prove. Otherwise $\Kd_a>0$ for every
$a\ne\as$; let $w_a:=1/\Kd_a\in[0,\infty)$. Let $\nu'\in\Alto(\nu;\mathcal E)$. Since
$\nu'_{\as}=\nu_{\as}$ and $\as$ is not optimal for $\nu'$, some arm $b\ne\as$ has
\begin{equation*}
  \marm_b(\nu')>\marm_{\as}(\nu')=\mean(\nu_{\as})=\ms.
\end{equation*}
As $\nu'_b\in\mathcal D_b$, the definition of $\Kinf$ gives $\KL(\nu_b,\nu'_b)\ge\Kd_b$; in
particular $\Kd_b<\infty$ (the set defining $\Kd_b$ contains $\nu'_b$), so $w_b=1/\Kd_b\in(0,\infty)$ and
\begin{equation*}
  \sum_{a\ne\as}w_a\KL(\nu_a,\nu'_a)\ge w_b\KL(\nu_b,\nu'_b)\ge\frac{\Kd_b}{\Kd_b}=1.
\end{equation*}
So $w$ is reduced-feasible, with cost $\sum_{a\ne\as}\gap_a/\Kd_a=S$, and $V\le S$.

Steps 2 and 3 give $V=S$, hence $\GL(\nu;\mathcal E)=S$. The case of a single model is
$\mathcal D_a=\mathcal D$ for every $a$.
\end{proof}

\subsection{Proof of Proposition~\ref{prop:gl-zero}}\label{sec:proof-gl-zero}

\begin{proof}
\uses{def:gl, lem:gl-reduction}
Let $\nu'\in\Alt(\nu;\mathcal E)$ with $\nu'_{\as}=\nu_{\as}$. By hypothesis $\nu'=\nu$, but $\as$ is
optimal for $\nu$ and not for $\nu'$, a contradiction. So $\Alto(\nu;\mathcal E)=\emptyset$, the
vector $w=0$ satisfies the constraints of Lemma~\ref{lem:gl-reduction} vacuously, and
$\GL(\nu;\mathcal E)\le0$; hence $\GL(\nu;\mathcal E)=0$.
\end{proof}

\subsection{Proof of Corollary~\ref{cor:lower}}\label{sec:proof-lower}

\begin{proof}
\uses{thm:main, def:consistent, def:gl, cite:GMS19-thm1, lem:classes, def:regret}
By Lemma~\ref{lem:classes}(iii), $\pi$ is uniformly fast convergent on $\powc{\mathcal D}$. Let $\as$ be
the optimal arm of $\nu$. For each $a\ne\as$, $\gap_a>0$, and Theorem~\ref{cite:GMS19-thm1}
(\emph{the lower bound of Garivier, M\'enard and Stoltz}) gives
\begin{equation*}
  \liminf_{T\to\infty}\frac{\Er[\Np_a(T)]}{\log T}\ge\frac{1}{\Kinf(\nu_a,\ms;\mathcal D)}.
\end{equation*}
Since the liminf of a finite sum of nonnegative sequences is at least the sum of their liminfs (in
$[0,\infty]$), and $\gap_{\as}=0$,
\begin{align*}
  \liminf_{T\to\infty}\frac{\Reg_T(\pi,\nu)}{\log T}
  &=\liminf_{T\to\infty}\sum_{a\ne\as}\gap_a\frac{\Er[\Np_a(T)]}{\log T}\\
  &\ge\sum_{a\ne\as}\frac{\gap_a}{\Kinf(\nu_a,\ms;\mathcal D)} && \text{(Theorem~\ref{cite:GMS19-thm1})}\\
  &=\GL(\nu;\powc{\mathcal D}) && \text{(Theorem~\ref{thm:main})}.\qedhere
\end{align*}
\end{proof}

% ----------------------------------------------------------------
\section{Proofs of the properties of the models and of \texorpdfstring{$\mathcal{K}_{\inf}$}{Kinf}}\label{sec:proof-properties}

In this appendix we prove the lemmas of Section~\ref{sec:properties}, in their order. Two of them use
the mixture of a law with a Dirac mass, so we first state that bound.

\begin{lemma}[Mixing with a point mass]\label{lem:mix}
\uses{def:laws}
For every law $P$, every $x\in\R$ and every $\eta\in[0,1)$, the law
$Q:=(1-\eta)P+\eta\delta_x$ satisfies $\KL(P,Q)\le\log\frac1{1-\eta}$. If
$P\in\Pone$, then $Q\in\Pone$ and $\mean(Q)=(1-\eta)\mean(P)+\eta x$.
\end{lemma}

\begin{proof}
\uses{def:laws}
Since $Q\ge(1-\eta)P$, every $Q$-null set is $P$-null, so $P\ll Q$ and the density satisfies
$\frac{dP}{dQ}\le\frac1{1-\eta}$ $Q$-almost everywhere. Hence
\begin{equation*}
  \KL(P,Q)=\int\log\frac{dP}{dQ}\,dP\le\log\frac1{1-\eta}.
\end{equation*}
The mean of $Q$ is linear in the mixture: $\int|y|\,dQ=(1-\eta)\int|y|\,dP+\eta|x|<\infty$
and $\mean(Q)=(1-\eta)\mean(P)+\eta x$.
\end{proof}

\begin{lemma}[The Bernoulli divergence]\label{lem:kl}
\uses{def:laws}
Let $p\in[0,1)$. The map $q\mapsto\kl(p,q)$ is finite, continuous and nondecreasing on $[p,1)$, with
$\kl(p,p)=0$, and $\kl(p,q)>0$ for $q\in(p,1]$.
\end{lemma}

\begin{proof}
\uses{def:laws}
For $q\in[p,1)$ with $q>0$, both terms of $\kl(p,q)$ are finite and smooth in $q$, and
\begin{equation*}
  \frac{\partial}{\partial q}\kl(p,q)=-\frac pq+\frac{1-p}{1-q}=\frac{q-p}{q(1-q)}\ge0.
\end{equation*}
If $p=0$, then $\kl(0,q)=\log\frac1{1-q}$, which is continuous and nondecreasing on $[0,1)$. In both
cases $q\mapsto\kl(p,q)$ is continuous and nondecreasing on $[p,1)$, and $\kl(p,p)=0$. For $q\in(p,1]$,
$\kl(p,q)=\KL(\Ber(p),\Ber(q))>0$ by Gibbs' inequality, since $\Ber(p)\ne\Ber(q)$.
\end{proof}

\subsection{Proof of Lemma~\ref{lem:kinf-basic}}\label{sec:proof-kinf-basic}

\begin{proof}
\uses{def:kinf}
Write $S(\mu,\mathcal D):=\{Q\in\mathcal D:\mean(Q)>\mu\}$, so that
$\Kinf(P,\mu;\mathcal D)=\inf_{Q\in S(\mu,\mathcal D)}\KL(P,Q)$.
\begin{enumerate}[label=(\roman*)]
  \item $S(\mu',\mathcal D)\subseteq S(\mu,\mathcal D)$ (since $\mu\le\mu'$), and an infimum over a
  smaller set is larger.
  \item $S(\mu,\mathcal D)\subseteq S(\mu,\mathcal D')$ (since $\mathcal D\subseteq\mathcal D'$), with
  the same reason.
  \item $S(\mu,\mathcal D)=\emptyset$, and $\inf\emptyset=+\infty$.
  \item $P\in S(\mu,\mathcal D)$, and $\KL(P,P)=0$, so the infimum is at most $0$; it is at least $0$
  since $\KL\ge0$.\qedhere
\end{enumerate}
\end{proof}

\subsection{Proof of Lemma~\ref{lem:kinf-unit}}\label{sec:proof-kinf-unit}

\begin{proof}
\uses{def:kinf, def:models, lem:kinf-basic, lem:mix, lem:kl, cite:GMS19-lem1}
\emph{Setting.} $P\in\Punit$, $\mu\in\R$. \emph{Goal.} Items (i)--(iii).

\emph{Item (i).} Every $Q\in\Punit$ has $\mean(Q)\le1\le\mu$, so no law of $\Punit$ has a mean larger
than $\mu$, and Lemma~\ref{lem:kinf-basic}(iii) (\emph{empty constraint set}) gives
$\Kinf(P,\mu;\Punit)=+\infty$.

\emph{Item (ii), lower bound.} Assume $\mean(P)\le\mu<1$ and let $Q\in\Punit$ with $\mean(Q)>\mu$.
We apply Lemma~\ref{cite:GMS19-lem1} (\emph{data processing for a $[0,1]$-valued variable}) on
$(\R,\mathcal B(\R))$ with $\mathbb P_1:=P$, $\mathbb P_2:=Q$ and $Z(x):=\min\{\max\{x,0\},1\}$. Since
$P$ and $Q$ are carried by $[0,1]$, $Z(x)=x$ almost surely under both, so $\E_P[Z]=\mean(P)$ and
$\E_Q[Z]=\mean(Q)$. Hence
\begin{align*}
  \KL(P,Q)&\ge\kl(\mean(P),\mean(Q)) && \text{(Lemma~\ref{cite:GMS19-lem1})}\\
          &\ge\kl(\mean(P),\mu) && \text{(Lemma~\ref{lem:kl}, since }\mean(P)\le\mu<\mean(Q)\le1\text{)}.
\end{align*}
For the second line, if $\mean(Q)<1$ the monotonicity of Lemma~\ref{lem:kl} on $[\mean(P),1)$
applies; if $\mean(Q)=1$, then $\kl(\mean(P),1)=+\infty$ because $\mean(P)\le\mu<1$. Taking the
infimum over $Q$ gives $\Kinf(P,\mu;\Punit)\ge\kl(\mean(P),\mu)$.

\emph{Item (ii), upper bound.} Let $\eta\in(\eta_0,1)$, where
$\eta_0:=\frac{\mu-\mean(P)}{1-\mean(P)}\in[0,1)$ (note $\mean(P)\le\mu<1$). Let
$Q_\eta:=(1-\eta)P+\eta\delta_1\in\Punit$. By Lemma~\ref{lem:mix} (\emph{mixing with a
point mass}),
\begin{equation*}
  \mean(Q_\eta)=\mean(P)+\eta(1-\mean(P))>\mean(P)+\eta_0(1-\mean(P))=\mu
  \qquad\text{and}\qquad\KL(P,Q_\eta)\le\log\frac1{1-\eta}.
\end{equation*}
So $\Kinf(P,\mu;\Punit)\le\log\frac1{1-\eta}$ for every $\eta\in(\eta_0,1)$, and
letting $\eta\downarrow\eta_0$,
\begin{equation*}
  \Kinf(P,\mu;\Punit)\le\log\frac1{1-\eta_0}=\log\frac{1-\mean(P)}{1-\mu}.
\end{equation*}

\emph{Item (ii), consequences.} If $\mean(P)<\mu<1$, the lower bound is positive by
Lemma~\ref{lem:kl} and the upper bound finite. If $\mu=\mean(P)<1$, the upper bound is $\log1=0$.

\emph{Item (iii).} Let $p\le\mu<1$. For $q\in(\mu,1]$, $\KL(\Ber(p),\Ber(q))=\kl(p,q)$, and the laws of
$\Bclass$ with a mean larger than $\mu$ are the $\Ber(q)$ with $q\in(\mu,1]$, so
\begin{equation*}
  \Kinf(\Ber(p),\mu;\Bclass)=\inf_{q\in(\mu,1]}\kl(p,q)=\kl(p,\mu),
\end{equation*}
because $q\mapsto\kl(p,q)$ is nondecreasing on $[p,1)$, continuous at $\mu$, and $\kl(p,1)=+\infty$
(Lemma~\ref{lem:kl}). Since $\Bclass\subseteq\Punit$, Lemma~\ref{lem:kinf-basic}(ii) and the lower
bound of item (ii) give
\begin{equation*}
  \kl(p,\mu)\le\Kinf(\Ber(p),\mu;\Punit)\le\Kinf(\Ber(p),\mu;\Bclass)=\kl(p,\mu).
\end{equation*}
Hence the three quantities are equal.
\end{proof}

\subsection{Proof of Lemma~\ref{lem:kinf-p1}}\label{sec:proof-kinf-p1}

\begin{proof}
\uses{def:kinf, lem:mix}
Let $P\in\Pone$, $\mu\in\R$ and $\eta\in(0,1)$. Let
\begin{equation*}
  x_\eta:=\frac{\mu+\eta-(1-\eta)\mean(P)}{\eta},\qquad
  Q_\eta:=(1-\eta)P+\eta\delta_{x_\eta}.
\end{equation*}
By Lemma~\ref{lem:mix} (\emph{mixing with a point mass}), $Q_\eta\in\Pone$,
$\mean(Q_\eta)=(1-\eta)\mean(P)+\eta x_\eta=\mu+\eta>\mu$ and
$\KL(P,Q_\eta)\le\log\frac1{1-\eta}$. Hence $0\le\Kinf(P,\mu;\Pone)\le\log\frac1{1-\eta}$
for every $\eta\in(0,1)$, and letting $\eta\to0$ gives $\Kinf(P,\mu;\Pone)=0$.
\end{proof}

\subsection{Proof of Lemma~\ref{lem:kinf-gauss}}\label{sec:proof-kinf-gauss}

\begin{proof}
\uses{def:kinf, def:models}
\emph{Setting.} $\mu_0\in\R$, $v>0$, $\mu\ge\mu_0$, $P:=\Nor(\mu_0,v)$. \emph{Goal.}
$\Kinf(P,\mu;\Gau)=\frac12\log(1+(\mu-\mu_0)^2/v)$.

For $Q=\Nor(\mu_1,v_1)$ with $\mu_1\in\R$ and $v_1>0$, the Gaussian formula gives
\begin{equation*}
  \KL(P,Q)=\frac12\log\frac{v_1}{v}+\frac{v+(\mu_0-\mu_1)^2}{2v_1}-\frac12.
\end{equation*}
Fix $\mu_1$ and write $d:=\mu_1-\mu_0$. As a function of $v_1>0$, the right-hand side has the
derivative $\frac1{2v_1}-\frac{v+d^2}{2v_1^2}$, which is negative for $v_1<v+d^2$ and positive for
$v_1>v+d^2$. So its minimum over $v_1>0$ is attained at $v_1=v+d^2$, where
\begin{equation*}
  \min_{v_1>0}\KL(\Nor(\mu_0,v),\Nor(\mu_1,v_1))=\frac12\log\frac{v+d^2}{v}=\frac12\log\Big(1+\frac{d^2}{v}\Big).
\end{equation*}
The laws of $\Gau$ with a mean larger than $\mu$ are the $\Nor(\mu_1,v_1)$ with $\mu_1>\mu$ and $v_1>0$.
Hence
\begin{equation*}
  \Kinf(P,\mu;\Gau)=\inf_{\mu_1>\mu}\frac12\log\Big(1+\frac{(\mu_1-\mu_0)^2}{v}\Big)
  =\frac12\log\Big(1+\frac{(\mu-\mu_0)^2}{v}\Big),
\end{equation*}
because $\mu_1-\mu_0>\mu-\mu_0\ge0$, the map $d\mapsto\frac12\log(1+d^2/v)$ is continuous and
increasing on $[0,\infty)$, and $\mu_1\downarrow\mu$ approaches the infimum.
\end{proof}

\subsection{Proof of Lemma~\ref{lem:semi}}\label{sec:proof-semi}

\begin{proof}
\uses{def:kinf, def:models, lem:kinf-basic, lem:mix, lem:kinf-unit}
\emph{Item (i).} If $\mu\ge1$, every $Q\in\Semi$ has $\mean(Q)\le1\le\mu$, and
Lemma~\ref{lem:kinf-basic}(iii) gives $+\infty$. If $\mu<1$ and $\mean(P)>\mu$, then
Lemma~\ref{lem:kinf-basic}(iv) gives $0<\infty$. If $\mean(P)\le\mu<1$, the argument of the upper
bound of Lemma~\ref{lem:kinf-unit}(ii) applies verbatim, since $Q_\eta=(1-\eta)P+\eta\delta_1\in\Semi$:
for every $\eta\in(\eta_0,1)$ with $\eta_0:=\frac{\mu-\mean(P)}{1-\mean(P)}$, Lemma~\ref{lem:mix}
gives $\mean(Q_\eta)>\mu$ and $\KL(P,Q_\eta)\le\log\frac1{1-\eta}$; letting
$\eta\downarrow\eta_0$ gives the bound $\log\frac{1-\mean(P)}{1-\mu}<\infty$. (Note
$\mean(P)\le\mu<1$, so $\eta_0\in[0,1)$.)

\emph{Item (ii).} A law $Q\in\Semi$ has $1-X\ge0$ $Q$-almost surely, so $\mean(Q)=1$ if and only if
$X=1$ $Q$-almost surely, that is $Q=\delta_1$. Hence $\ms=1$ if and only if some arm has the law
$\delta_1$.

\emph{Item (iii).} $\Punit\subseteq\Semi$, and Lemma~\ref{lem:kinf-basic}(ii) applies.

\emph{Item (iv).} For $Q\in\Lmom_{\varepsilon,B}$, Jensen's inequality for the convex map
$y\mapsto y^{1+\varepsilon}$ on $[0,\infty)$ gives
\begin{equation*}
  \Big(\int|x|\,dQ(x)\Big)^{1+\varepsilon}\le\int|x|^{1+\varepsilon}\,dQ(x)\le B,
\end{equation*}
hence $\int|x|\,dQ\le B^{1/(1+\varepsilon)}<\infty$. The law $\delta_0$ has
$\int|x|^{1+\varepsilon}\,d\delta_0=0\le B$, and $\delta_x$ has $\int|y|^{1+\varepsilon}\,d\delta_x=|x|^{1+\varepsilon}$.
\end{proof}

\subsection{Proof of Lemma~\ref{lem:kurt}}\label{sec:proof-kurt}

\begin{proof}
\uses{def:models, def:kinf}
\emph{Item (i).} Let $Y:=X-\mean(P)$ under $P$. By Jensen's inequality (or Cauchy--Schwarz),
$\E[Y^4]\ge(\E[Y^2])^2=\var(P)^2$, and dividing by $\var(P)^2>0$ gives $\kurt(P)\ge1$.

\emph{Item (ii).} Under $\phi_\#P$, the variable is $c+bX$, with mean $c+b\mean(P)$, centred
variable $bY$, variance $b^2\var(P)$ and centred fourth moment $b^4\E[Y^4]$. So
$\kurt(\phi_\#P)=b^4\E[Y^4]/(b^2\var(P))^2=\kurt(P)$, and the moment conditions of $\Hk_\kappa$ are
preserved by $\phi$ and by $\phi^{-1}(y)=(y-c)/b$. Hence $Q\mapsto\phi_\#Q$ is a bijection of
$\Hk_\kappa$ onto itself. Since $b>0$, $\mean(\phi_\#Q)>\phi(\mu)$ if and only if $\mean(Q)>\mu$; and
$\KL(\phi_\#P,\phi_\#Q)=\KL(P,Q)$ by invariance under the measurable bijection $\phi$. Taking the
infimum over $Q$ gives the equality of the two divergences.

\emph{Item (iii).} $\Ber(1/2)$ has $\var=1/4$ and centred fourth moment $1/16$, so its kurtosis is
$1\le\kappa$. A Gaussian law has kurtosis $3$, so it lies in $\Hk_\kappa$ if and only if $\kappa\ge3$.

\emph{Item (iv).} $\var(\delta_x)=0$. For $0<p<1$, $\Ber(p)$ has $\var=p(1-p)=:v$ and centred fourth
moment $p(1-p)^4+(1-p)p^4=v(1-3v)$, so
\begin{equation*}
  \kurt(\Ber(p))=\frac{v(1-3v)}{v^2}=\frac1v-3,
\end{equation*}
which is larger than $\kappa$ exactly when $v<1/(\kappa+3)$.
\end{proof}

% ----------------------------------------------------------------
\section{Proofs for the classes of learners}\label{sec:proof-learners}

In this appendix we prove the results of Section~\ref{sec:learners}: the relations between the
classes, the lemma on rewards by count, the bound on \ucb{} (by a decomposition lemma and one lemma
per term), and the lemma on constant and uniform learners.

\subsection{Proof of Lemma~\ref{lem:classes}}\label{sec:proof-classes}

\begin{proof}
\uses{def:consistent, def:ufc, def:smarter, def:regret}
\emph{Item (i).} Each property is a statement for every $\nu$ of the class; it holds for every
$\nu\in\mathcal E$ when it holds for every $\nu\in\mathcal E'$.

\emph{Item (ii).} Every $p>\gamma'$ satisfies $p>\gamma$.

\emph{Item (iii).} Let $\nu\in\mathcal E$, $a$ a suboptimal arm and $p\in(0,1]$. Since every term of
$\Reg_T(\pi,\nu)=\sum_b\gap_b\Er[\Np_b(T)]$ is nonnegative,
\begin{equation*}
  \frac{\Er[\Np_a(T)]}{T^p}\le\frac{1}{\gap_a}\,\frac{\Reg_T(\pi,\nu)}{T^p}\xrightarrow[T\to\infty]{}0.
\end{equation*}

\emph{Item (iv).} Let $\nu\in\mathcal E$ and $p>0$. For $T\ge2$, the arms with $\gap_b=0$ do not
contribute to the regret, so
\begin{equation*}
  \frac{\Reg_T(\pi,\nu)}{T^p}=\sum_{b:\gap_b>0}\frac{\gap_b\Er[\Np_b(T)]}{T^p}
  \le\sum_{b:\gap_b>0}\frac{C}{\gap_b}\,\frac{\log T}{T^p}\xrightarrow[T\to\infty]{}0.
\end{equation*}
So $\pi$ is consistent on $\mathcal E$, and uniformly fast convergent by item (iii).

\emph{Item (v).} The regret of $(\pi)_{T\ge1}$ at horizon $T$ is $\Reg_T(\pi,\nu)$.
\end{proof}

\subsection{Proof of Lemma~\ref{lem:stack}}\label{sec:proof-stack}

\begin{proof}
\uses{def:run, def:regret}
\emph{Idea.} We replace, in the array of rewards by count, the rewards collected at rounds $u\ge t$
by the independent variables $\Zc$. At $t=0$ this array is $\Zc$, with the product law. From $t$ to
$t+1$ a single coordinate changes, and it has the conditional law $\nu_{\Ap_t}$ before and after.
Letting $t\to\infty$ gives the array $\Yc$.

\emph{Setting.} Write $I:=\arms\times\mathbb N_{\ge1}$ and $\mu_I:=\bigotimes_{(a,s)\in I}\nu_a$, the
law of $\Zc=(\Zc_i)_{i\in I}$ on $\R^I$. \emph{Goal.} $\Yc=(\Yc_i)_{i\in I}$ has the law $\mu_I$, and
$\hm_a(t)=\bY_{a,\Np_a(t)}$.

\emph{Step 1: the truncated arrays.} For $t\ge0$ and $i\in I$, let $\Yc^t_i:=\Xr_{\taup_i}$ if
$\taup_i<t$, and $\Yc^t_i:=\Zc_i$ otherwise. Since $\taup_i<t$ only if $\taup_i\le t-1$, the event
$\{\taup_i<t\}$ and the value $\Xr_{\taup_i}$ on it are functions of $\Hist_t$. So
$\Yc^t=\Phi_t(\Hist_t,\Zc)$ for a measurable map $\Phi_t$, in which the coordinate $i$ of the
output depends on $\Zc$ only through $\Zc_i$. At $t=0$ no $\taup_i$ is below $0$, so $\Yc^0=\Zc$
has the law $\mu_I$.

\emph{Step 2: the law of $\Yc^t$ does not depend on $t$.} Fix $t\ge0$ and let
$W:=(\Hist_t,\Ap_t)$ and $c:=(\Ap_t,\Np_{\Ap_t}(t)+1)\in I$, the index of the pull of round $t$, so
that $\taup_c=t$; $c$ is a function $c(W)$ of $W$. For $i\ne c$, $\taup_i\ne t$, so
$\{\taup_i<t+1\}=\{\taup_i<t\}$ and $\Yc^{t+1}_i=\Yc^t_i$. For $i=c$, $\Yc^{t+1}_c=\Xr_t$ and
$\Yc^t_c=\Zc_c$. Hence there is a measurable map $G$ such that, for every bounded measurable
$F:\R^I\to\R$,
\begin{equation*}
  F(\Yc^{t+1})=G(W,\Xr_t,\Zc),\qquad F(\Yc^t)=G(W,\Zc_{c(W)},\Zc),
\end{equation*}
where $G(w,x,z)$ is $F$ applied to the array $\Phi_t(h,z)$ (for $w=(h,a)$) with its coordinate $c(w)$
replaced by $x$; in particular $G(w,x,z)$ does not depend on $z_{c(w)}$. We compute both
expectations.

First, $\Zc$ is independent of $(W,\Xr_t)$, and the conditional law of $\Xr_t$ given $W$ is
$\nu_{\Ap_t}$ (Definition~\ref{def:run}); so the conditional law of $\Xr_t$ given $(W,\Zc)$ is
$\nu_{\Ap_t}$, and
\begin{equation*}
  \E[F(\Yc^{t+1})]=\E\Big[\int G(W,x,\Zc)\,\nu_{\Ap_t}(dx)\Big].
\end{equation*}
Second, $W$ is independent of $\Zc$. For a fixed $w=(h,a)$, under $\mu_I$ the coordinate $z_{c(w)}$
has the law $\nu_a$ and is independent of the other coordinates, on which alone $G(w,x,z)$ depends
besides $x$; by Fubini's theorem,
\begin{equation*}
  \int G(w,z_{c(w)},z)\,\mu_I(dz)=\int\!\!\int G(w,x,z)\,\nu_a(dx)\,\mu_I(dz).
\end{equation*}
Integrating in $w$ against the law of $W$ (Fubini again, by independence) gives
\begin{equation*}
  \E[F(\Yc^t)]=\E\Big[\int G(W,x,\Zc)\,\nu_{\Ap_t}(dx)\Big]=\E[F(\Yc^{t+1})].
\end{equation*}
So $\Yc^{t+1}$ and $\Yc^t$ have the same law, and by induction from Step 1 every $\Yc^t$ has the law
$\mu_I$.

\emph{Step 3: the limit.} For each $i\in I$, $\Yc^t_i=\Yc_i$ for every $t>\taup_i$ if $\taup_i<\infty$,
and for every $t$ if $\taup_i=\infty$. So $\Yc^t\to\Yc$ coordinatewise. For a bounded continuous
$F$ that depends on finitely many coordinates, dominated convergence gives
$\E[F(\Yc)]=\lim_t\E[F(\Yc^t)]=\int F\,d\mu_I$. Such functions determine the finite-dimensional laws,
and the finite-dimensional laws determine a law on $\R^I$; hence $\Yc$ has the law $\mu_I$, that is
the $\Yc_{a,s}$ are independent with $\Yc_{a,s}\sim\nu_a$.

\emph{Step 4: empirical means.} Let $t\ge0$ and $a$ with $s:=\Np_a(t)\ge1$. The rounds $u<t$ with
$\Ap_u=a$ are exactly $\taup_{a,1}<\dots<\taup_{a,s}$, all finite and below $t$, so
\begin{equation*}
  \hm_a(t)=\frac1s\sum_{u=0}^{t-1}\Xr_u\ind(\Ap_u=a)=\frac1s\sum_{j=1}^s\Xr_{\taup_{a,j}}=\frac1s\sum_{j=1}^s\Yc_{a,j}=\bY_{a,s}.\qedhere
\end{equation*}
\end{proof}

\subsection{Proof of Proposition~\ref{prop:ucb}}\label{sec:proof-ucb}

Throughout this subsection, $\nu\in\powc{\Punit}$, $a$ is a suboptimal arm, $\gap:=\gap_a$, $\as$ is
an optimal arm of $\nu$ (any of them), $T\ge1$, and we work on the enlarged run of
Lemma~\ref{lem:stack}. To prove the proposition we will use the following decomposition. Let
\begin{equation*}
  u:=\max\Big\{1,\Big\lceil\frac{8\log T}{\gap^2}\Big\rceil\Big\},
\end{equation*}
and, for $1\le t\le T-1$, let $E_t$ be the event that the optimal arm is underestimated:
\begin{equation*}
  E_t:=\big\{\Np_{\as}(t)\ge1\big\}\cap\big\{\Uix_{\as}(t)\le\ms\big\}.
\end{equation*}

\begin{lemma}[Decomposition]\label{lem:ucb-decomp}
\uses{def:ucb, lem:stack}
Almost surely,
\begin{equation*}
  \Np_a(T)\le u+\underbrace{\sum_{t=1}^{T-1}\ind(E_t)}_{\mathsf A}
  +\underbrace{\sum_{s=u}^{\infty}\ind\Big(\bY_{a,s}-\marm_a>\frac{\gap}{2}\Big)}_{\mathsf B}.
\end{equation*}
\end{lemma}
(see the proof in Subsection~\ref{sec:proof-ucb-decomp}).

From this decomposition we will use the following lemma to control the two random terms.

\begin{lemma}[The two terms]\label{lem:ucb-terms}
\uses{lem:stack, def:ucb, def:models}
With the notation of Lemma~\ref{lem:ucb-decomp},
\begin{equation*}
  \E[\mathsf A]\le\frac14,\qquad\E[\mathsf B]\le\frac{2}{\gap^2}.
\end{equation*}
\end{lemma}
(see the proof in Subsection~\ref{sec:proof-ucb-terms}). Both bounds are classical, from
Hoeffding's inequality; no contribution here.

\begin{proof}[Proof of Proposition~\ref{prop:ucb}]
\proves{prop:ucb}
\uses{lem:ucb-decomp, lem:ucb-terms, lem:classes, def:ucb, def:ufc, def:consistent, def:models}
By Lemmas~\ref{lem:ucb-decomp} and~\ref{lem:ucb-terms}, and since $\lceil y\rceil\le y+1$ for $y\ge0$,
\begin{align*}
  \Er[\Np_a(T)]&\le u+\E[\mathsf A]+\E[\mathsf B]\\
  &\le1+\frac{8\log T}{\gap^2}+\frac14+\frac{2}{\gap^2},
\end{align*}
which is the first bound. The means lie in $[0,1]$, so $\gap\le1$ and $\frac54\le\frac{5}{4\gap^2}$;
for $T\ge2$, $\frac{13}{4}\le5\log2\le5\log T$. Hence, for $T\ge2$,
\begin{equation*}
  \Er[\Np_a(T)]\le\frac{8\log T+\frac{13}{4}}{\gap^2}\le\frac{13\log T}{\gap^2}.
\end{equation*}
This holds for every instance of $\powc{\Punit}$, every suboptimal arm and every tie rule, which is
uniform super-fast convergence on $\powc{\Punit}$ with the constant $13$. Lemma~\ref{lem:classes}(iv)
gives consistency and uniform fast convergence on $\powc{\Punit}$, and Lemma~\ref{lem:classes}(i), with
$\powc{\Bclass}\subseteq\powc{\Punit}$, gives them on $\powc{\Bclass}$.
\end{proof}

\subsubsection{Proof of Lemma~\ref{lem:ucb-decomp}}\label{sec:proof-ucb-decomp}

\begin{proof}
\uses{def:ucb, lem:stack}
Almost surely, $\Ap_t\in M_t(\Hist_t)$ for every $t$, since
$\Pru(\Ap_t\in M_t(\Hist_t))=\E[\rho_t(\Hist_t,M_t(\Hist_t))]=1$. We split the rounds at which $a$ is
pulled according to its count:
\begin{equation*}
  \Np_a(T)=\sum_{t=0}^{T-1}\ind(\Ap_t=a,\ \Np_a(t)<u)+\sum_{t=0}^{T-1}\ind(\Ap_t=a,\ \Np_a(t)\ge u).
\end{equation*}
Let us look at the first sum. The count $\Np_a(t)$ increases by one at each pull of $a$, so the
rounds $t$ with $\Ap_t=a$ have distinct values of $\Np_a(t)$; those counted have
$\Np_a(t)\in\{0,\dots,u-1\}$, so the first sum is at most $u$.

Now the second sum. Let $t\le T-1$ with $\Ap_t=a$ and $s:=\Np_a(t)\ge u\ge1$. Then $t\ge1$ and
$\Uix_a(t)<\infty$. If $\Np_{\as}(t)=0$, then $\Uix_{\as}(t)=+\infty>\Uix_a(t)$, which contradicts
$a\in M_t(\Hist_t)$; so $\Np_{\as}(t)\ge1$. Then we split according to $E_t$:
\begin{equation*}
  \ind(\Ap_t=a,\ \Np_a(t)=s)\le\ind(E_t)+\ind\big(\Ap_t=a,\ \Np_a(t)=s,\ \Uix_{\as}(t)>\ms\big).
\end{equation*}
On the event of the last indicator, $a$ is a maximiser of the index, so
\begin{align*}
  \ms&<\Uix_{\as}(t)\le\Uix_a(t) && (a\in M_t(\Hist_t))\\
  &=\bY_{a,s}+\sqrt{\frac{2\log(t+1)}{s}} && \text{(Lemma~\ref{lem:stack}: }\hm_a(t)=\bY_{a,\Np_a(t)}\text{)}\\
  &\le\bY_{a,s}+\sqrt{\frac{2\log T}{u}} && (t+1\le T,\ s\ge u)\\
  &\le\bY_{a,s}+\frac{\gap}{2} && \Big(u\ge\frac{8\log T}{\gap^2}\Big).
\end{align*}
Since $\ms-\gap/2=\marm_a+\gap/2$, this event is contained in $\{\Ap_t=a,\ \Np_a(t)=s,\ \bY_{a,s}-\marm_a>\gap/2\}$.
For each $s$ at most one round $t$ has $\Ap_t=a$ and $\Np_a(t)=s$, so summing over $t$ and $s\ge u$,
\begin{equation*}
  \sum_{t=0}^{T-1}\ind(\Ap_t=a,\ \Np_a(t)\ge u)\le\sum_{t=1}^{T-1}\ind(E_t)+\sum_{s=u}^{\infty}\ind\Big(\bY_{a,s}-\marm_a>\frac\gap2\Big).
\end{equation*}
Adding the two sums gives the decomposition.
\end{proof}

\subsubsection{Proof of Lemma~\ref{lem:ucb-terms}}\label{sec:proof-ucb-terms}

\begin{proof}
\uses{lem:stack, def:ucb, def:models}
By Lemma~\ref{lem:stack} (\emph{rewards by count}), for each arm $b$ the variables $\Yc_{b,1},\Yc_{b,2},\dots$
are independent with the law $\nu_b\in\Punit$, so Hoeffding's inequality gives, for every $s\ge1$ and
$x\ge0$,
\begin{equation}\label{eq:hoeffding}
  \Pru\big(\bY_{b,s}-\marm_b\le-x\big)\le e^{-2sx^2},\qquad\Pru\big(\bY_{b,s}-\marm_b\ge x\big)\le e^{-2sx^2}.
\end{equation}

\emph{Term $\mathsf A$.} Let $1\le t\le T-1$. On $E_t$, $s:=\Np_{\as}(t)\in\{1,\dots,t\}$ and, by
Lemma~\ref{lem:stack}, $\Uix_{\as}(t)=\bY_{\as,s}+\sqrt{2\log(t+1)/s}$. A union bound over the
possible values of $s$ and \eqref{eq:hoeffding} with $x=\sqrt{2\log(t+1)/s}$ give
\begin{align*}
  \Pru(E_t)&\le\sum_{s=1}^{t}\Pru\Big(\bY_{\as,s}-\ms\le-\sqrt{\frac{2\log(t+1)}{s}}\Big)
  \le\sum_{s=1}^{t}e^{-4\log(t+1)}=\frac{t}{(t+1)^4}\le\frac{1}{(t+1)^3}.
\end{align*}
Hence
\begin{equation*}
  \E[\mathsf A]\le\sum_{t=1}^{\infty}\frac{1}{(t+1)^3}=\zeta(3)-1\le\frac14.
\end{equation*}

\emph{Term $\mathsf B$.} By \eqref{eq:hoeffding} with $x=\gap/2$, and since $e^y-1\ge y$,
\begin{equation*}
  \E[\mathsf B]\le\sum_{s=1}^{\infty}e^{-s\gap^2/2}=\frac{1}{e^{\gap^2/2}-1}\le\frac{2}{\gap^2}.
\end{equation*}
(The event of \eqref{eq:hoeffding} with $\ge$ contains the event with $>$.)
\end{proof}

\subsection{Proof of Lemma~\ref{lem:constant}}\label{sec:proof-constant}

\begin{proof}
\uses{def:consistent, def:ufc, def:smarter, def:regret, def:models}
\emph{Item (i).} Under $\Cons^b$, $\Ap_t=b$ almost surely for every $t$, so $\Np_b(T)=T$, $\Np_a(T)=0$
for $a\ne b$, and $\Reg_T=\gap_bT$.

\emph{Item (ii).} Under $\pi^{\mathrm{u}}$, the conditional law of $\Ap_t$ given $\Hist_t$ is uniform,
so $\Pru(\Ap_t=a)=1/\Karms$ and $\E[\Np_a(T)]=T/\Karms$; the regret follows from
Definition~\ref{def:regret}.

\emph{Consequences.} Let $\nu\in\mathcal E$ with $\gap_b>0$. For $p\in(0,1)$, $\Reg_T(\Cons^b,\nu)/T^p=
\gap_bT^{1-p}\to\infty$, so $\Cons^b$ is not $\gamma$-consistent for any $\gamma\in[0,1)$ (take
$p\in(\gamma,1)$); $\Er[\Np_b(T)]/T^{1}=1\not\to0$, so it is not uniformly fast convergent; it is
not uniformly super-fast convergent with any constant $C$, since $T>C\log T/\gap_b^2$ for $T$ large;
and an optimal arm $a\ne b$ of $\nu$ has $\Er[\Np_a(T)]=0<T/\Karms$, so it is not smarter than the
uniform learner. The class $\powc{\Bclass}$ contains such an instance: give $b$ the law $\Ber(0)$ and
every other arm the law $\Ber(1)$. For $\pi^{\mathrm{u}}$, $T/\Karms\ge T/\Karms$ holds for every arm,
hence for the optimal ones; and if $\gap_a>0$ for some $a$, then $\Reg_T(\pi^{\mathrm{u}},\nu)/T^p\ge
\gap_aT^{1-p}/\Karms\to\infty$ for $p\in(\gamma,1)$.
\end{proof}

\subsection{Proof of Lemma~\ref{lem:gms-classes}}\label{sec:proof-gms-classes}

\begin{proof}
\uses{def:pairwise, def:monotonic, def:models, def:learner, def:run, lem:constant}
\emph{Item (i).} Let $\nu$ be an instance and $T\ge1$. Under $\pi^{\mathrm{u}}$, the arms
$\Ap_0,\dots,\Ap_{T-1}$ are independent and uniform on $\arms$: by Definition~\ref{def:run}, for every
$t<T$ and $a_0,\dots,a_t\in\arms$,
\begin{align*}
  \Pru(\Ap_0=a_0,\dots,\Ap_t=a_t)&=\E\big[\ind(\Ap_0=a_0,\dots,\Ap_{t-1}=a_{t-1})\,\pol_t(\Hist_t,\{a_t\})\big]\\
  &=\frac1{\Karms}\Pru(\Ap_0=a_0,\dots,\Ap_{t-1}=a_{t-1}),
\end{align*}
and induction on $t$ gives $\Karms^{-t-1}$. Let $a,a'$ be two arms and $\sigma$ the transposition of
$a$ and $a'$ on $\arms$. The uniform law on $\arms^T$ is invariant under $\sigma$ applied to each
coordinate, and this map exchanges the numbers of coordinates equal to $a$ and to $a'$. Hence
$(\Np_a(T),\Np_{a'}(T))$ and $(\Np_{a'}(T),\Np_a(T))$ have the same law, for every pair of arms, in
particular for optimal arms with equal laws. For monotonicity, let $\nu'\preccurlyeq\nu$; then
$\Opt(\nu')=\Opt(\nu)$, and by Lemma~\ref{lem:constant}(ii) both sides of the inequality of
Definition~\ref{def:monotonic} are $|\Opt(\nu)|\,T/\Karms$.

\emph{Item (ii).} Let $\nu$ be such an instance and $T\ge1$. Under $\Cons^b$, $(\Np_b(T),\Np_a(T))=(T,0)$
almost surely, while $(\Np_a(T),\Np_b(T))=(0,T)$; the two laws differ. The instance of
$\powc{\Bclass}$ in which every arm has the law $\Ber(1/2)$ is such an instance, since $\Karms\ge2$.

\emph{Item (iii).} The learner $\pi^{\mathrm{sw}}$ is given by the kernels $\pol_0:=\delta_b$ and, for
$t\ge1$ and $h=(a_0,x_0,\dots)\in\Hs_t$, $\pol_t(h,\cdot):=\delta_b$ if $x_0\le1/2$ and $\delta_c$
otherwise, which are measurable in $h$. Let $\nu\in\powc{\Bclass}$ have $\nu_b=\Ber(1/2)$ and
$\nu_a=\Ber(1)$ for $a\ne b$, and let $\nu'$ be $\nu$ with $\nu'_b:=\Ber(0)$. Then
$\Opt(\nu)=\arms\setminus\{b\}$, the optimal arms keep their laws and $\marm_b(\nu')=0\le\frac12$, so
$\nu'\preccurlyeq\nu$. Under $\nu$, $\Ap_0=b$ and $\Xr_0$ has the law $\Ber(1/2)$; for $t\ge1$,
$\Ap_t=c$ if $\Xr_0=1$ and $\Ap_t=b$ otherwise. So $\sum_{a\in\Opt(\nu)}\Np_a(2)=\Np_c(2)=\ind(\Xr_0=1)$,
with expectation $\frac12$. Under $\nu'$, $\Xr_0=0$ almost surely, so $\Ap_t=b$ for every $t$ and the
sum is $0<\frac12$. The inequality of Definition~\ref{def:monotonic} fails at $T=2$.
\end{proof}

\subsection{Proof of Lemma~\ref{lem:tails}}\label{sec:proof-tails}

\begin{proof}
\uses{def:tails, def:regret, def:models, lem:constant}
\emph{Item (i).} Let $\nu\in\mathcal E$ have a unique optimal arm $\as$, and $n\ge2$. If
$\Np_a(n)<C\log n/\gap_a^2$ for every $a\ne\as$, then, since $\gap_{\as}=0$ and $\gap_a\ge\gap_{\min}$,
\begin{equation*}
  \hR_n=\sum_{a\ne\as}\gap_a\Np_a(n)<\sum_{a\ne\as}\frac{C\log n}{\gap_a}\le\frac{(\Karms-1)C\log n}{\gap_{\min}}.
\end{equation*}
So the event $\{\hR_n\ge(\Karms-1)C\log n/\gap_{\min}\}$ is contained in the union over the $\Karms-1$
suboptimal arms $a$ of $\{\Np_a(n)\ge C\log n/\gap_a^2\}$, whose probability is at most
$(\Karms-1)\tilde C/f(n)$.

\emph{Item (ii).} Let $\nu\in\mathcal E$ have a unique optimal arm, with $\gap_b>0$, and let
$C,\tilde C>0$. Under $\Cons^b$, $\Np_b(n)=n$ and $\hR_n=n\gap_b$ almost surely (Lemma~\ref{lem:constant}(i)).
There is $n_0$ such that $n\ge C\log n/\gap_b^2$ and $n\gap_b\ge C\log n/\gap_{\min}$ for every $n\ge n_0$;
for these $n$ both probabilities of Definition~\ref{def:tails} equal $1$. Since $f$ is unbounded, it
is not bounded by $\tilde C$ on $\{n_0,n_0+1,\dots\}$ (it would then be bounded), so some $n\ge n_0$ has
$\tilde C/f(n)<1$, and neither inequality holds for the constants $C,\tilde C$. The instance of
$\powc{\Bclass}$ with $\nu_c=\Ber(1)$ for one arm $c\ne b$ and $\nu_a=\Ber(0)$ for the others has the
unique optimal arm $c$, and $b$ is suboptimal in it.
\end{proof}

\subsection{Proof of Proposition~\ref{prop:ucb-tails}}\label{sec:proof-ucb-tails}

\begin{proof}
\uses{def:ucb, def:tails, def:models, lem:stack, lem:tails}
\emph{Setting.} $\nu\in\powc{\Punit}$, $a$ a suboptimal arm, $\gap:=\gap_a\in(0,1]$, $\as$ an optimal
arm, $T\ge2$; we work on the enlarged run of Lemma~\ref{lem:stack}, and $E_t$ is the event
$\{\Np_{\as}(t)\ge1,\ \Uix_{\as}(t)\le\ms\}$ of Subsection~\ref{sec:proof-ucb}. \emph{Goal.}
$\Pru(\Np_a(T)\ge x)\le(\log T)^{-2}$, where $x:=10\log T/\gap^2$.

\emph{Step 1: a single count.} Let $u:=\lceil x\rceil$ and $s:=u-1$. Since $\gap\le1$ and
$2\log2>1$,
\begin{equation*}
  s\ge x-1=\frac{8\log T}{\gap^2}+\frac{2\log T}{\gap^2}-1>\frac{8\log T}{\gap^2}\ge8\log2>5 .
\end{equation*}
As $\Np_a$ is an integer, $\{\Np_a(T)\ge x\}=\{\Np_a(T)\ge u\}$. On this event, since $\Np_a(t)$ increases
by one at each pull of $a$, some round $t\le T-1$ has $\Ap_t=a$ and $\Np_a(t)=s$; then $t\ge s$, since
$\Np_a(t)\le t$.

\emph{Step 2: the event at that round.} Almost surely $\Ap_t\in M_t(\Hist_t)$ for every $t$ (as in
the proof of Lemma~\ref{lem:ucb-decomp}). Let $t$ be as in Step 1. Since $s\ge1$, $\Uix_a(t)<\infty$, so
$\Np_{\as}(t)\ge1$ (otherwise $\Uix_{\as}(t)=+\infty>\Uix_a(t)$). If $E_t$ fails, then, by
Lemma~\ref{lem:stack} and $t+1\le T$,
\begin{equation*}
  \ms<\Uix_{\as}(t)\le\Uix_a(t)=\bY_{a,s}+\sqrt{\frac{2\log(t+1)}{s}}\le\bY_{a,s}+\sqrt{\frac{2\log T}{s}}
  <\bY_{a,s}+\frac\gap2,
\end{equation*}
that is $\bY_{a,s}-\marm_a>\gap/2$. Hence
\begin{equation*}
  \{\Np_a(T)\ge x\}\subseteq\Big\{\bY_{a,s}-\marm_a>\frac\gap2\Big\}\cup\bigcup_{t=s}^{T-1}E_t.
\end{equation*}

\emph{Step 3: the probabilities.} By Lemma~\ref{lem:stack}, the $\Yc_{b,j}$ are independent with laws
$\nu_b\in\Punit$, and Hoeffding's inequality gives
$\Pru(\bY_{a,s}-\marm_a>\gap/2)\le e^{-s\gap^2/2}\le e^{-4\log T}=T^{-4}$. For $t\ge1$, on $E_t$ the
count $j:=\Np_{\as}(t)$ lies in $\{1,\dots,t\}$ and $\bY_{\as,j}-\ms\le-\sqrt{2\log(t+1)/j}$, so a
union bound over $j$ and Hoeffding's inequality give $\Pru(E_t)\le t\,(t+1)^{-4}\le(t+1)^{-3}$. Since
$(t+1)^{-3}\le\int_t^{t+1}y^{-3}\,dy$,
\begin{equation*}
  \sum_{t=s}^{T-1}\Pru(E_t)\le\int_s^\infty y^{-3}\,dy=\frac1{2s^2}\le\frac{\gap^4}{128(\log T)^2}\le\frac1{128(\log T)^2}.
\end{equation*}
Finally $(\log T)^2\le T^2\le T^4/2$ for $T\ge2$, so $T^{-4}\le\frac12(\log T)^{-2}$, and
\begin{equation*}
  \Pru\Big(\Np_a(T)\ge\frac{10\log T}{\gap^2}\Big)\le\frac{1}{2(\log T)^2}+\frac{1}{128(\log T)^2}\le\frac{1}{(\log T)^2}.
\end{equation*}
This is Definition~\ref{def:tails} for the counts, with $f(n)=(\log n)^2$, $C=10$ and $\tilde C=1$, on the
instances of $\powc{\Punit}$ with a unique optimal arm; Lemma~\ref{lem:tails}(i) gives the pseudo-regret.
\end{proof}

\subsection{Proof of Lemma~\ref{lem:family}}\label{sec:proof-family}

\begin{proof}
\uses{def:family}
When $\pi^T=\pi$ for every $T$, the law and the expectations of $\pi^T$ at each horizon are those of
$\pi$, so the condition of Definition~\ref{def:family} for the family is, word for word, the condition
of the corresponding definition for $\pi$.
\end{proof}

% ----------------------------------------------------------------
\section{Sources of the techniques}\label{sec:attribution}

\begin{center}
\small
\begin{tabular}{@{}p{0.47\textwidth}p{0.45\textwidth}@{}}
\toprule
Technique (where) & Earliest source found \\
\midrule
Data processing for a $[0,1]$-valued variable (Lemma~\ref{lem:kinf-unit}(ii)) & \citet[Lemma~1]{GMS19}; the resulting bound $\Kinf\ge\kl$ on $\Punit$ is in \citet[Section~6.1]{CGMMS13} \\
Mixing with a Dirac mass to raise the mean (Lemmas~\ref{lem:mix}, \ref{lem:kinf-unit}, \ref{lem:kinf-p1}, \ref{lem:semi}) & source not known; see also the mixing step of \citet[Appendix~B.1]{GHMS22}, the dual formula of \citet[Theorem~8]{HT10} (the same upper bound by Jensen's inequality) and, for $\Kinf=0$, \citet[Exercise~16.4]{LS} \\
Weight $\infty$ on the optimal arm, alternatives keeping its law (Lemma~\ref{lem:gl-reduction}) & the constraint set of \citet[Theorem~1]{CMP17}; see also \citet{ATA89,GL97} \\
Alternatives moving one arm (Theorem~\ref{thm:main}, Step 1) & \citet{LR85}, \citet[proof of Theorem~1]{BK96}; in this form, the proof of \citet[Theorem~1]{GMS19} \\
Feasibility of the rates $1/\Kinf$ through one arm above $\ms$ (Theorem~\ref{thm:main}, Step 3) & \citet[Section~4]{CMP17} (parametric), \citet[Example~3]{VPG24} (finite set of rewards) \\
Truncation in time for rewards by count (Lemma~\ref{lem:stack}) & the formalisation in \citet{LML}; see also \citet{D36} and \citet[Section~4.6]{LS} \\
Decomposition and Hoeffding bounds for \ucb{} (Propositions~\ref{prop:ucb} and~\ref{prop:ucb-tails}) & \citet[Theorem~1]{ACF02} \\
Gaussian closed form of $\Kinf$ (Lemma~\ref{lem:kinf-gauss}) & \citet[Section~4.4]{BK96} \\
Uniform, constant and switching learners as exclusions (Lemmas~\ref{lem:constant}, \ref{lem:gms-classes}, \ref{lem:tails}) & elementary; source not known \\
\bottomrule
\end{tabular}
\end{center}

% ----------------------------------------------------------------
\section{Notation}\label{sec:notation}

Every symbol is a link to its definition.
\begin{center}
\small
\begin{tabular}{@{}lll@{}}
\toprule
Symbol & Meaning & Defined in \\
\midrule
$\Pall$, $\Pone$, $\mean(P)$ & laws, laws with a finite mean, mean & Definition~\ref{def:laws} \\
$\KL$, $\kl$, $\Ber(p)$ & divergences, Bernoulli law & Definition~\ref{def:laws} \\
$\Kinf(P,\mu;\mathcal D)$ & divergence to the laws of $\mathcal D$ with a larger mean & Definition~\ref{def:kinf} \\
$\Punit$, $\Bclass$, $\Semi$, $\Lmom_{\varepsilon,B}$ & models & Definition~\ref{def:models} \\
$\var(P)$, $\kurt(P)$, $\Hk_\kappa$ & variance, kurtosis, kurtosis model & Definition~\ref{def:models} \\
$\Nor(\mu,v)$, $\Gau$ & Gaussian law, Gaussian model & Definition~\ref{def:models} \\
$\arms$, $\Karms$ & arms, number of arms & Definition~\ref{def:instance} \\
$\marm_a$, $\ms$, $\gap_a$, $\Opt(\nu)$ & means, best mean, gaps, optimal arms & Definition~\ref{def:instance} \\
$\prodc{\mathcal D}$, $\powc{\mathcal D}$ & unstructured classes & Definition~\ref{def:product} \\
$\Hs_t$, $\pol_t$ & histories, kernels of a learner & Definition~\ref{def:learner} \\
$\Ap_t$, $\Xr_t$, $\Hist_t$ & arm, reward, history of round $t$ & Definition~\ref{def:run} \\
$\Pru_{\pi,\nu}$, $\Er_{\pi,\nu}$ & law and expectation of a run & Section~\ref{sec:setting}, after Definition~\ref{def:run} \\
$\Np_a(t)$, $\hm_a(t)$ & count and empirical mean before round $t$ & Definition~\ref{def:regret} \\
$\hR_T$, $\Reg_T(\pi,\nu)$ & pseudo-regret, regret & Definition~\ref{def:regret} \\
$\Zc_{a,s}$, $\taup_{a,s}$, $\Yc_{a,s}$, $\bY_{a,s}$ & auxiliary array, pull times, rewards by count & Lemma~\ref{lem:stack} \\
$\Uix_a(t)$ & index of \ucb{} & \eqref{eq:ucb} \\
$\Cons^b$, $\pi^{\mathrm{u}}$ & constant and uniform learners & Lemma~\ref{lem:constant} \\
$\nu'\preccurlyeq\nu$ & the other arms of $\nu'$ are worse & before Definition~\ref{def:monotonic} \\
$\gap_{\min}$ & smallest gap & before Definition~\ref{def:tails} \\
$\pi^{\mathrm{sw}}$ & switching learner & Lemma~\ref{lem:gms-classes} \\
$\as$, $\Alt(\nu;\mathcal E)$ & optimal arm, alternatives & Definition~\ref{def:gl} \\
$\GL(\nu;\mathcal E)$ & Graves--Lai constant & Definition~\ref{def:gl} \\
$\Alto(\nu;\mathcal E)$ & alternatives keeping the optimal law & Lemma~\ref{lem:gl-reduction} \\
$\Kd_a$ & $\Kinf(\nu_a,\ms;\mathcal D_a)$ in the proof of Theorem~\ref{thm:main} & Appendix~\ref{sec:proof-main} \\
\bottomrule
\end{tabular}
\end{center}

Defined words: \term{w:cons}{consistent} and \term{w:acons}{$\gamma$-consistent}
(Definition~\ref{def:consistent}), \term{w:ufc}{uniformly fast convergent} and
\term{w:usfc}{uniformly super-fast convergent} (Definition~\ref{def:ufc}),
\term{w:smart}{smarter than the uniform learner} (Definition~\ref{def:smarter}),
\term{w:pairwise}{pairwise symmetric for optimal arms} (Definition~\ref{def:pairwise}),
\term{w:monotonic}{monotonic} (Definition~\ref{def:monotonic}), \term{w:tailT}{$f$-upper tailed
counts} and \term{w:tailR}{pseudo-regret} (Definition~\ref{def:tails}), and their
\term{w:family}{horizon-dependent versions} (Definition~\ref{def:family}).

\bibliographystyle{plainnat}
\bibliography{refs,../../papers/library}

\end{document}
